
% Introduction --------------------------------

% \section{Introduction}
% \label{sec:intro}

% Many real-world spatio-temporal systems, such as environmental processes, fluid transport, and other simulated physical phenomena, are governed by partial differential equations (PDEs) that encode conservation laws and structural dynamics. A general form of such systems is 
% \begin{tightalign}
% \partial_t \mathbf{u}(\mathbf{x}, t) 
% &= \mathcal{F}\big(
%     \mathbf{u}(\mathbf{x}, t),\ 
%     \nabla \mathbf{u}(\mathbf{x}, t),\ 
%     \Delta \mathbf{u}(\mathbf{x}, t),\ 
%     \ldots;\ 
%     \boldsymbol{\theta} 
% \big), \nonumber \\
% &\quad \mathbf{x} \in \Omega \subset \mathbb{R}^d,\quad t \in [0, T]
% \end{tightalign}
% where \( \mathbf{u}(\mathbf{x}, t) \in \mathbb{R}^C \) represents the evolving physical state (e.g., concentration or velocity fields), and \( \boldsymbol{\theta} \) denotes the latent parameters of the system, such as diffusivity, advection velocity, or source intensity.

% In practice, observations are limited to sparse and noisy sensor measurements, typically obtained at \( N_s \) spatial locations \( \{\mathbf{s}_i\}_{i=1}^{N_s} \), which may be misaligned with the discretized simulation grid. The data consist of time-indexed sensor traces
% \begin{tightalign}
% \mathbf{d}_t^{(i)} = \mathcal{O}(\mathbf{u}_t; \mathbf{s}_i) + \boldsymbol{\varepsilon}_t^{(i)}, \quad 
% \boldsymbol{\varepsilon}_t^{(i)} \sim \mathcal{N}(\mathbf{0}, \boldsymbol{\Sigma}_i),
% \end{tightalign}
% where \( \mathcal{O} \) is a differentiable observation operator (e.g., bilinear sampling), and \( \mathbf{d}_t \in \mathbb{R}^{N_s \times m} \) represents partial and noisy views of the underlying state. Here, \( m \leq C \) denotes the number of state components observed by each sensor, allowing for partial observations of the full physical state.

% This leads to the central learning problem: recover the governing parameters \( \boldsymbol{\theta} \) such that the simulated dynamics \( \Phi_{\boldsymbol{\theta}}(\mathbf{u}_0, t) \) match the observed sensor behavior. Formally, we seek to solve
% \begin{tightalign}
% \min_{\boldsymbol{\theta}} \sum_{t=1}^T \mathcal{L}_{\text{obs}}\left( \mathcal{O}(\mathbf{u}_t),\; \mathbf{d}_t \right) + \lambda\, \mathcal{R}(\boldsymbol{\theta}),
% \end{tightalign}
% where \( \mathbf{u}_t = \Phi_{\boldsymbol{\theta}}(\mathbf{u}_0, t) \) is the simulated state at time \( t \), \( \mathcal{O}(\mathbf{u}_t) \in \mathbb{R}^{N_s \times m} \) denotes the observation operator applied to the state, \( \mathcal{L}_{\text{obs}} \) is a data fidelity loss (e.g., masked or weighted squared error), and \( \mathcal{R}(\boldsymbol{\theta}) \) encodes structural priors or regularization over the parameter field.

% This inverse problem is fundamentally underdetermined due to sparse and noisy supervision. Multiple parameter configurations can match limited sensor data, and additional challenges such as misalignment, missing values, and heteroscedastic noise further destabilize learning. Classical inverse solvers often rely on strong initialization or dense supervision to remain stable. Modern machine learning approaches such as PINNs~\cite{raissi2019physics} and neural operators~\cite{li2024physics,li2020fourier,rahman2022u} encode physics but degrade under sparse observations. In contrast, purely statistical models can capture sensor-level dynamics but lack physical consistency. Bridging this gap requires a framework that integrates physical structure with data-driven generalization, enabling robust learning under real-world sensing constraints.


% We introduce \textbf{STAMP}, a physics-guided learning framework that combines a differentiable simulator \( \Phi_{\boldsymbol{\theta}} \) with a frozen spatio-temporal prior to recover latent physical parameters from sparse observations. Unlike physics-informed networks that require dense supervision or statistical models that ignore governing equations, STAMP couples the structure of a finite-difference PDE solver with the flexibility of a learned sensor-space model. The prior is implemented as a message-passing recurrent network trained separately on sensor traces, then held fixed during parameter optimization. It regularizes learning by encouraging consistency with empirical dynamics when physical gradients are underconstrained. STAMP jointly minimizes a composite loss over simulation error \( \mathcal{L}_{\text{obs}} \), prior alignment, and regularization on \( \boldsymbol{\theta} \), enabling accurate and generalizable inference in low-observability regimes.

% This paper makes the following contributions:
% \begin{enumerate}[topsep=0pt, leftmargin=*, itemsep=0pt]
%     \item \textbf{Hybrid learning framework:} We propose a method that integrates a differentiable PDE-constrained simulator with a learned spatio-temporal prior, enabling robust model calibration from sparse and noisy observations.
%     \item \textbf{Structured simulator design:} We develop a fully differentiable finite-difference solver with constrained parameterization and boundary-aware updates, supporting stable gradient-based optimization of physical parameters.
%     \item \textbf{Frozen message-passing prior:} We introduce a frozen recurrent graph-based prior trained on sensor trajectories, which regularizes dynamics learning under limited observability.
%     \item \textbf{Generalizable training objective:} We formulate a joint loss that balances physical fidelity, empirical trajectory consistency, and observational alignment, leading to improved generalization across varying sensor configurations and data regimes.
% \end{enumerate}

% STAMP achieves accurate recovery of spatially varying physical parameters and reproduces observed dynamics under extreme sparsity and sensor noise. It outperforms both physics-only and purely data-driven baselines, demonstrating improved calibration, physical consistency, and generalization to unseen boundary and forcing conditions. \todo{Empirical result summary here.} More broadly, STAMP provides a general strategy for learning structured physical models in underdetermined regimes by combining differentiable simulation with learned sensor-space priors, contributing to the development of physically grounded and data-efficient machine learning systems.





% Methods ------------------------------------------------

% \section{Method}

% We introduce STAMP, a hybrid learning framework that combines a differentiable finite-difference (FD) simulator with a frozen spatio-temporal message-passing prior. This section describes the physical simulator, observation model, prior dynamics network, and training objective.
% \todo{Add a concept figure on top of the page.}

% \paragraph{Physical Model and Simulator}
% We consider systems governed by PDEs of the form
% \begin{tightalign}
% \partial_t \mathbf{u}(\mathbf{x}, t) = \mathcal{F}(\mathbf{u}; \boldsymbol{\theta}),
% \end{tightalign}
% where \( \mathbf{u}_t \in \mathbb{R}^{H \times W \times C} \) denotes the discretized physical state over space and time, and \( \boldsymbol{\theta} \) are unknown parameters.

% \noindent\textbf{\textit{Finite-Difference Core}} We define a differentiable simulator \( \Phi_{\boldsymbol{\theta}} \) such that \( \mathbf{u}_{t+1} = \Phi_{\boldsymbol{\theta}}(\mathbf{u}_t) \) using explicit or semi-implicit FD schemes. Spatial operators (e.g., gradients, divergence, Laplacian) are implemented with convolutions or roll-based stencils, supporting upwind or Rusanov fluxes for stability. Semi-implicit updates (e.g., Jacobi iterations) are used for stiff terms. Parameters \( \boldsymbol{\theta} \) may be scalars or spatial fields with constrained reparameterizations (e.g., softplus) and smoothness priors (TV, Laplacian). Boundary conditions (periodic, Dirichlet, Neumann) are enforced via ghost cells or padding policies.
% \todo{Revise and add details after implementation.}

% \noindent\textbf{\textit{Observation and Sensor Model}}
% We observe the system through \( N_s \) fixed spatial sensors at locations \( \{ \mathbf{s}_i \}_{i=1}^{N_s} \). The observation at time \( t \) is given by
% \begin{tightalign}
% \hat{\mathbf{d}}_t^{(i)} = \mathcal{O}(\mathbf{u}_t; \mathbf{s}_i + \boldsymbol{\delta}_i),
% \end{tightalign}
% where \( \boldsymbol{\delta}_i \) is a learnable offset with prior \( \mathcal{N}(\mathbf{0}, \sigma_s^2 \mathbf{I}) \), and \( \mathcal{O} \) is bilinear sampling. Observed data \( \mathbf{d}_t \in \mathbb{R}^{N_s \times m} \) may include missing values (masked by \( \mathbf{M}_t \)) and heteroscedastic noise with known variances \( \{\sigma_i^2\} \).
% \todo{Revise and add details after implementation.}

% \paragraph{Spatio-Temporal Dynamics Prior} We train a message-passing recurrent neural network (MPRNN) on ground-truth sensor traces \( \{ \mathbf{d}_t \} \) to model empirical dynamics. The model is trained with a next-step prediction loss and early stopping on held-out traces. Once trained, its parameters are frozen and used as a nonparametric teacher prior \( f_{\text{MPRNN}} \) to guide physical model calibration under limited supervision.

% \noindent\textbf{\textit{MPRNN Architecture}}
% MPRNN extends the message-passing neural network (MPNN) framework \citep{gilmer2017message} to spatio-temporal forecasting on sensor graphs. The network maintains a hidden state \( h_{v,t} \) for each sensor node \( v \), updated over time through a combination of local recurrence and spatial message passing. At each time step, the hidden state is first updated via an observation function \( O(h_{v,t-1}, d_{v,t-1}) \), followed by aggregation of messages from spatial neighbors:
% \begin{tightalign*}
%   m_{v,t+1} &= \sum_{w \in N(v)} M(h_{v,t}, h_{w,t}, e_{v,w}) \\
%   h_{v,t+1} &= U(h_{v,t}, m_{v,t+1}) \\
%   \hat{d}_{v,t+1} &= R(h_{v,t+1}),
% \end{tightalign*}
% where \( M \) and \( U \) are learnable message and update functions, \( R \) is a decoder, and \( e_{v,w} \) denotes edge features such as spatial proximity. This recurrence allows the model to learn spatio-temporal patterns from partial and noisy sensor trajectories. In settings requiring broader spatial context, the update can include multiple rounds of intra-time message passing before the temporal transition. 


% % \paragraph{Training Objective}
% % STAMP optimizes a composite objective:
% % \begin{itemize}
% %   \item \textbf{Data loss (masked, weighted):}
% %   \begin{tightalign}
% %   \mathcal{L}_{\text{data}} = \sum_t \left\| \mathbf{M}_t \odot (\hat{\mathbf{d}}_t - \mathbf{d}_t) \right\|_2^2.
% %   \end{tightalign}
% %   \item \textbf{Dynamics consistency loss:}
% %   \begin{tightalign}
% %   \mathcal{L}_{\text{dyn}}^{(2)} = \sum_{t=1}^{T-1} \left\| f_{\text{MPRNN}}(\hat{\mathbf{d}}_t) - f_{\text{MPRNN}}(\mathbf{d}_t) \right\|_2^2.
% %   \end{tightalign}
% %   \item \textbf{Physics priors:}
% %   \begin{tightalign}
% %   \mathcal{L}_{\text{phys}} = \lambda_{\text{sm}} \|\nabla \boldsymbol{\theta}\|_2^2 + \lambda_{\text{con}}\,\text{Drift}(\mathbf{u}_{1:T}) + \lambda_{\text{rng}}\,\text{Barrier}(\boldsymbol{\theta}).
% %   \end{tightalign}
% % \end{itemize}
% % The total loss is:
% % \begin{tightalign}
% % \mathcal{L} = \lambda_d \mathcal{L}_{\text{data}} + \lambda_y \mathcal{L}_{\text{dyn}} + \lambda_p \mathcal{L}_{\text{phys}}.
% % \end{tightalign}

% \paragraph{Training Objective}
% STAMP jointly optimizes physical consistency, empirical regularization, and observation fidelity through a composite loss function:

% \noindent\textbf{\textit{Data loss (masked, weighted):}} This term encourages agreement between the simulated state and the observed sensor traces. Sensor-specific noise levels and missing data are handled via heteroscedastic weights and binary masks:
% \begin{tightalign}
% \mathcal{L}_{\text{data}} = \sum_t \left\| \mathbf{M}_t \odot \left(\hat{\mathbf{d}}_t - \mathbf{d}_t\right) \right\|_2^2,
% \end{tightalign}
% where \( \hat{\mathbf{d}}_t = \mathcal{O}(\mathbf{u}_t) \) is the observation of the simulated state at time \( t \), \( \mathbf{d}_t \) is the corresponding ground truth sensor measurement, and \( \mathbf{M}_t \in \{0,1\}^{N_s \times m} \) masks out missing entries.

% \noindent\textbf{\textit{Dynamics consistency loss:}} This term enforces alignment between the empirical dynamics encoded by the frozen MPRNN teacher and the simulated observations. It promotes consistency in the latent trajectory space even when direct data supervision is weak:
% \begin{tightalign}
% \mathcal{L}_{\text{dyn}}^{(2)} = \sum_{t=1}^{T-1} \left\| f_{\text{MPRNN}}(\hat{\mathbf{d}}_t) - f_{\text{MPRNN}}(\mathbf{d}_t) \right\|_2^2.
% \end{tightalign}
% This form of teacher forcing regularizes learning in underdetermined regions by encouraging the simulator to produce trajectories that conform to empirical spatio-temporal structure.

% \noindent\textbf{\textit{Physics priors:}} To ensure plausible physical parameters, we include a regularization term composed of multiple priors:
% \begin{tightalign}
% \mathcal{L}_{\text{phys}} = \lambda_{\text{sm}} \|\nabla \boldsymbol{\theta}\|_2^2 
% + \lambda_{\text{con}}\,\text{Drift}(\mathbf{u}_{1:T}) 
% + \lambda_{\text{rng}}\,\text{Barrier}(\boldsymbol{\theta}).
% \end{tightalign}
% Here, \( \nabla \boldsymbol{\theta} \) penalizes spatial roughness in the parameter field, \( \text{Drift}(\mathbf{u}_{1:T}) \) penalizes non-conservation of mass or energy over time, and \( \text{Barrier}(\boldsymbol{\theta}) \) enforces range constraints to keep parameters within physically meaningful bounds.
% \todo{Revise and add new losses after implementation.}

% \noindent The total objective combines these terms with tunable weights:
% \begin{tightalign}
% \mathcal{L} = \lambda_d \mathcal{L}_{\text{data}} + \lambda_y \mathcal{L}_{\text{dyn}} + \lambda_p \mathcal{L}_{\text{phys}}.
% \end{tightalign}
% Each term plays a complementary role: the data loss anchors the model to observed measurements, the dynamics loss enforces learned temporal structure, and the physics priors encode inductive constraints for regularity and conservation.


% \paragraph{Optimization and Training}
% We use truncated backpropagation-through-time (BPTT) with window size \( K \), sliding across episodes and detaching hidden states. Training stabilizes with curriculum over \( K \), adaptive \( \Delta t \), and gradient clipping. Mini-batches sample diverse ICs/BCs to avoid degenerate solutions. Sensor offsets \( \boldsymbol{\delta}_i \) may be learned jointly with \( \boldsymbol{\theta} \).









% Theory --------------------------------------

% \section{Theoretical Foundations}
% \label{sec:theory}
% In this section, we provide the theoretical analysis of MPRNN and some of its properties.

% \subsection{Eqivalence between MPRNN and Spatio-Temporal Processes}

% A spatio-temporal process refers to a phenomenon that is dependent on both the spatial and temporal coordinates. Generally, such a process is defined as $P:(\Vec{s},t) \rightarrow \mathbb{R}$, where, the spatial encoding is represented as a vector in some coordinate system, and the time is represented by a real number, and the output value, another real number, is the target of interest. Most climate processes can be generally expressed as spatio-temporal processes.

% Let us first formalize our notion of spatio-temporal processes.
% \begin{definition}
% A spatio-temporal process is a process defined on 4-dimensional Euclidean space-time coordinate system x,y,z,t by partial differential equations of the following form.
% \begin{align*}
%     \frac{\partial^{m} u}{\partial t^{m}} = i(t) + a\frac{\partial^{n} u}{\partial x^{n}} + b\frac{\partial^{n} u}{\partial y^{n}} + c\frac{\partial^{n} u}{\partial z^{n}}
% \end{align*}
% We also assume knowledge of initial values $u(x,y,z,0)$, $u^{(1)}(x,y,z,0)$, ..., $u^{(m-1)}(x,y,z,0)$.
% \end{definition}
% In actual examples of climate processes, we limit the value of $m \in \{1,2\}$. For such examples of climate processes, we note the following result.
% \begin{theorem}
% Any spatio-temporal process defined between the interval $D_{xyzt} = [(0,0,0,0), (X,Y,Z,T)]$, with given initial conditions $u(x,y,z,0), u^{(1)}(x,y,z,0), \dots,$ $ u^{(m-1)}(x,y,z,0), \forall (x,y,z) \in D_{xyz}$ can be approximated using message passing recurrent neural networks with arbitrarily small error.
% \label{thm:mprnn_spatiotemporal}
% \end{theorem}
% The proof of this result can be found in Appendix \ref{sec:proof1}. It implies that the MPRNN architecture is has the mathematical complexity to represent a large category of climate processes. This expressive power is key to scaling MPRNNs to a large variety of climate processes, and creating system models that simultaneously model a large number of climate parameters.

% \subsection{Convergence of MPRNNs}
% From the results proved by \cite{hardt2018gradient}, we know that when RNNs are trained to learn time invariant and controllable linear systems with Gaussian error using stochastic gradient descent, the learned parameters $\hat{\Theta}$ have a population risk
% \begin{align*}
%     f(\hat{\Theta}) \leq f(\Theta) + O(\sqrt{\frac{n^5+\sigma^2n^3}{TN}})
% \end{align*}
% when given N sample sequences of length T. Here, $f(\Theta)$ is the optimal population risk, $n$ is the order of the system (dimension of the hidden state), $\sigma^2$ is the variance of the Gaussian error, and the population risk is defined as 
% \begin{align}
%     \label{eqn:population_risk}
%     f(\hat{\Theta}) = \mathbb{E}[\frac{1}{T}\sum_{t=1}^T||\hat{u}_t-u_t||^2]
% \end{align}
% where $u_t$ and $\hat{u}_t$ are the actual and predicted outputs of the dynamical system.

% Extending this result to MPRNN training on spatio-temporal processes, we note the following result.
% \begin{theorem}
% \label{thm:mprnn_convergence}
% Under the same assumptions as \cite{hardt2018gradient}, each individual node of MPRNN can learn parameters $\hat{\Theta}$ with the population risk
% \begin{align*}
%     f(\hat{\Theta}) \leq f(\Theta) + O(\sqrt{\frac{n^5+n^3V(u)(\Delta x \Delta t + (\Delta t)^2)^2}{TN}})
% \end{align*}
% assuming a uniform spatial grid ($\Delta x = \Delta y = \Delta z$), given N sample sequences of length T. Here, V(u) is an equation specific function.
% \end{theorem}
% The proof of this result is presented in Appendix \ref{sec:proof2}. This implies that the error of training the MPRNN architecture depends on the discretization resolution of the PDE solver ($\Delta t, \Delta x$), as well as the amount of training data ($T, N$). Discounting the effects of different PDEs (since we want to support all types of PDEs), we encounter a trade-off between accuracy of our model and computational complexity (higher resolution, more training samples).






% Evaluation --------------------------

% \section{Evaluation}
% \label{sec:evaluation}

% \subsection{Baselines and Metrics}

% \subsection{Evaluation Settings}
% \todo{Placeholders copied from FieldFormer for now. Update and revise after implementation.}

% \noindent\textbf{\textit{Periodic Heat Equation (2D)}}
% We first evaluate on a controlled synthetic setting governed by the anisotropic two-dimensional heat equation with periodic boundary conditions,
% \[
% u_t = \alpha_x u_{xx} + \alpha_y u_{yy} + f(x,y,t),
% \]
% where the forcing term is a smooth sinusoidal function of space and time. The domain is discretized on a $64 \times 64$ spatial grid with 10,000 time steps, and stability is enforced via the CFL condition. A smooth sinusoidal initial condition is evolved forward using explicit finite differences with periodic wrapping. From the resulting full field, we sample 20 sensor locations uniformly at random and extract both clean and noisy time series. Noise is scaled relative to the maximum magnitude, ensuring realistic observation corruption. The dataset is stored in compressed format with complete metadata (grid, parameters, sensor indices) for reproducibility. This benchmark provides a clean, scalar yet anisotropic setting where all physical parameters and boundary behavior are fully known. 

% \noindent\textbf{Periodic Shallow Water Equations (2D)}
% We next evaluate on a more complex synthetic setting governed by the linearized two-dimensional shallow water equations with periodic boundary conditions,
% \[
% \eta_t + H(u_x + v_y) = 0, \qquad 
% u_t + g \eta_x = 0, \qquad 
% v_t + g \eta_y = 0,
% \]
% which describe the propagation of free-surface \textit{gravity waves}. Unlike the purely diffusive heat equation, this system is hyperbolic and supports oscillatory, wave-like solutions with characteristic speed $c = \sqrt{gH}$. A Gaussian bump in surface height is used as the initial condition, generating outward-traveling gravity waves. The domain is discretized on a $64 \times 64$ spatial grid with 10,000 time steps, and numerical stability is enforced through a conservative Courant--Friedrichs--Lewy (CFL) condition for the forward--backward time stepping scheme. From the resulting fields $(\eta,u,v)$, we sample 20 sensor locations to obtain clean and noisy time series. 

% \noindent\textbf{\textit{Real-World Advection--Diffusion Pollution Simulation}} 
% To evaluate our methods on a realistic transport--diffusion regime, we simulate pollutant dispersion over New Delhi using the two-dimensional advection--diffusion equation,
% \[
% u_t + \mathbf{v}(t)\cdot\nabla u \;=\; \kappa \nabla^2 u + S(x,y),
% \]
% on a normalized $[0,1]\times[0,1]$ domain corresponding to the city’s geographic extent. The initial concentration field is constructed from the 32 regulatory monitoring stations\citep{cpcb_portal} on June 1, 2018, the first day of monsoon, interpolated onto the model grid using Universal Kriging. The source field combines emissions from brick kilns and industries sourced from \cite{GUTTIKUNDA2013101}, population density from \cite{ColumbiaUniversity2018}, and traffic intensity layers from Google Maps, cropped to the simulation window and normalized. 
% Advection is driven by a synthetic \emph{monsoon-mode wind} oriented toward the northeast, scaled from a 3.5~m/s base flow, based on averaged wind-speeds data from \cite{cisl_rda_ds083.2}, and time-compressed so that five simulation seconds correspond to one real day. To capture realistic variability, the wind includes diurnal strengthening and weakening, a small directional wobble, and autoregressive (AR(1)) stochastic gusts with multi-hour correlation times. Horizontal turbulent mixing is represented by a diffusivity chosen to yield an advection-dominated Péclet number at the city scale. Open lateral boundaries are treated with Orlanski-type radiation conditions supplemented by a sponge layer to minimize reflections. 

% We discretize the PDE with finite differences---upwind for advection, central for diffusion---and advance with a two-stage Heun (RK2) scheme. The simulation produces full spatio-temporal concentration fields as well as time series at sensor gridpoints, to which observation noise is added to reflect sensor variability.

% \subsection{Results}

% \subsection{Ablations}






% Appendix ----------------------

% \section{Proofs}
% \subsection{Proof of Theorem \ref{thm:mprnn_spatiotemporal}}
% \label{sec:proof1}
% \begin{proof}
% Let us first consider the case with n=1. We can write the corresponding spatio-temporal process as
% \begin{align*}
%     \frac{\partial^{m} u}{\partial t^{m}} = i(t) + a.\frac{\partial u}{\partial x} + b.\frac{\partial u}{\partial y} + c.\frac{\partial u}{\partial z}
% \end{align*}
% Say we look at one spatial point $p_0 = (x_0,y_0,z_0)$. The process can be written as
% \begin{align*}
%     \frac{d^{m} u}{d t^{m}}|_{p_0} = i(t) + a.\frac{\partial u}{\partial x}|_{p_0} + b.\frac{\partial u}{\partial y}|_{p_0} + c.\frac{\partial u}{\partial z}|_{p_0}
% \end{align*} 
% Assuming that the values of $\frac{\partial u}{\partial x}|_{p_0}, \frac{\partial u}{\partial y}|_{p_0}, \frac{\partial u}{\partial z}|_{p_0}$ is known, the point can be modelled by an RNN arbitrarily closely with error $O(\Delta t)$, based on lemma \ref{lemma:temporal_dynamical_equivalence}.

% Assume that we have discretized the spatial domain into a grid, with step sizes $\Delta x, \Delta y, \Delta z$. Then, 
% \begin{align*}
%     \frac{\partial u}{\partial x}|_{p_0} = \frac{u_{x_0+\Delta x, y_0, z_0, t} - u_{x_0, y_0, z_0,t}}{\Delta x} + O(\Delta x)
% \end{align*}
% based on Taylor series expansion. Thus, for this m-order point-temporal process, the error in the input term can be approximated by $O(\Delta x + \Delta y + \Delta z)$. But, if we have already approximated $u_{x_0+\Delta x, y_0, z_0, t}$ and $u_{x_0, y_0, z_0, t}$ using RNNs, there is another corresponding approximation error of $O(\Delta t)$ (based on lemma \ref{lemma:temporal_dynamical_equivalence}). Thus, the final error in the input is of $O(\Delta x + \Delta y + \Delta z + \Delta t)$.

% Then, we can rewrite the process equation as
% \begin{align*}
%     \frac{d^{m} u}{d t^{m}}|_{p_0} = i'(t) + O(\Delta x + \Delta y + \Delta z + \Delta t)
% \end{align*}
% where,
% \begin{align*}
%     i'(t) = i(t) + a.\frac{\partial u}{\partial x}|_{p_0} + b.\frac{\partial u}{\partial y}|_{p_0} + c.\frac{\partial u}{\partial z}|_{p_0}
% \end{align*}
% Naturally, the corresponding error for the discrete computation for next timestep will be
% \begin{align*}
%     u_{p_0,t+\Delta t} &= u_{p_0,t} + (i'(t) +\\ &O(\Delta x + \Delta y + \Delta z + \Delta t))*\Delta t + O(\Delta t^2)
% \end{align*}
% \begin{align*}
%     \implies u_{p_0,t+\Delta t} &= u_{p_0,t} + i'(t)*\Delta t +\\ & O((\Delta x + \Delta y + \Delta z)\Delta t + \Delta t^2)
% \end{align*}
% The compounded global error will be
% \begin{align*}
%     E = O(\Delta x + \Delta y + \Delta z + \Delta t)
% \end{align*}
% Now, we show how the input  $\frac{\partial u}{\partial x}|_{p_0}, \frac{\partial u}{\partial y}|_{p_0}, \frac{\partial u}{\partial z}|_{p_0}$ is obtained using message passing in the MPRNN framework. Assume that we form a grid graph between the discretized point in the domain $D_{xyz}$. Then, each point will carry a state and an output value. Let the readout function $R(h) = h$, making the state equal to the output. Further, let the message function $M(h_{v,t}, h_{w,t}, e_{v,w}) = h_{w,t}$, which means that each node sends its own state as its message. Then, the aggregation function to generate the message does a weighted sum of all messages along the x,y, and z directions.
% \begin{align*}
%     m_{v,t+1} = \frac{a}{\Delta x}h_{v+\Delta x, t}+ \frac{b}{\Delta y}h_{v+\Delta y, t}+ \frac{c}{\Delta z}h_{v+\Delta z, t}
% \end{align*}
% Also, the update function $U_v(h_{v,t}, m_{v,t+1})$, can be written as
% \begin{align*}
%     U(h_{v,t}, m_{v,t+1}) &= h_{v,t} + (i(t) + m_{v,t+1}\\ &- \frac{h_{v,t}}{\Delta x} - \frac{h_{v,t}}{\Delta y} - \frac{h_{v,t}}{\Delta z})*\Delta t
% \end{align*}
% This formulation shows how the messages in the MPRNN framework can produce the input $i'(t)$. Combined with the previous result, it shows that MPRNN can approximate spatio-temporal processes with $n=1$ to arbitrary accuracy under sufficient discretization. Now, for extending to cases $n>1$, we note that
% \begin{align*}
%     \frac{\partial^n u}{\partial x^n}|_{p_0} = \frac{\Delta^n_{\Delta x} u}{(\Delta x)^n} + O(\Delta x)
% \end{align*}
% based on \cite{wikipedia_finite_difference}. Thus, the error obtained using the finite difference remains the same, keeping the calculations unchanged. It should be noted that for computing nth order differences for approximating the nth order differential equations, at least $\frac{n}{2}$ message-passing steps have to be done before each update step. Also, based on lemma \ref{lemma:temporal_dynamical_equivalence}, we can generalize for m.
% \end{proof}

% \begin{lemma}
% \label{lemma:temporal_dynamical_equivalence}
% An nth order (point-)temporal process can be approximated by an n-stacked RNNs with arbitrary accuracy.
% \end{lemma}

% \begin{proof}
% First, we note that RNNs are universal approximators for open dynamical systems based on \cite{10.1007/11840817_66}. Thus, our proof reduces to showing that any nth order (point-)temporal process can be approximated by n sequentially connected open dynamical systems with arbitrary accuracy for time interval [0,T]. 

% First we consider the base case of first order (point-)temporal process. From Taylor series expansion, we have
% \begin{align*}
%     u(t+\delta t) =  u(t) + i(t)*\delta t + \frac{\frac{d^2u}{dt^2}(t)*(\delta t)^2}{2!} + ... 
% \end{align*}
% Then, we can rewrite the above as:
% \begin{align*}
%     u(t+\delta t) = u(t) + i(t)*\delta t + E
% \end{align*}
% where,
% \begin{align*}
%     E = O(\delta t^2) \text{ for } \delta t < 1
% \end{align*}
% Now, if we discretize the above differential equation as an open dynamical system such that
% \begin{align*}
%     h_{t+1} &= h_t + i(t)*\Delta t\\
%     u_t  &= h_t
% \end{align*}
% with the initial conditions $h_0 = u_0$, and step size $\Delta t$. Then the error between the outputs of the discrete and continuous formulations at time T will be
% \begin{align*}
%     E(T) = O(\Delta t^2)*(T)/\Delta t = O(\Delta t * T) = O(\Delta t)
% \end{align*}
% Thus, for any arbitrary error bound, one can find a small step size such that E(T) is within bounds.

% Now, assume that an $k^{th}$ order point-temporal process can be approximated by k connected open dynamical systems. We can write the $k+1^{th}$ order point-temporal process as the system of differential equations of the form
% \begin{align*}
%     \frac{d^{k+1}u}{dt^{k+1}} = \frac{du^{(k)}}{dt}\\
%     \frac{du^{(k)}}{dt} = i(t)
% \end{align*}
% Again, using Taylor series expansion
% \begin{align*}
%     u^{(k+1)}(t+\delta t) =  u^{(k+1)}(t) + \frac{du^{(k)}}{dt}*\delta t + \\
%     \frac{\frac{d^2u^{(k+1)}}{dt^2}(t)*(\delta t)^2}{2!} + ... \\
%     \implies u^{(k+1)}(t+\delta t) =  u^{(k+1)}(t) + \frac{du^{(k)}}{dt}*\delta t + E'
% \end{align*}
% Now, let's assume that we replace the $\frac{du^{(k)}}{dt}$ term in the above equation by the connected dynamical system estimate. Then, as $u^{(k)}_t = \frac{du^{(k)}}{dt} + O(\Delta t)$, we have,
% \begin{align*}
%     u^{(k+1)}(t+\delta t) =  u^{(k+1)}(t) + (u^{(k)}_t + O(\Delta t))*\delta t + E'
% \end{align*}
% where $E' = O(\delta t^2)$. Now, as we estimate the above discretization using a dynamical system, we will have
% \begin{align*}
%     u^{(k+1)}(t+\Delta t) &=  u^{(k+1)}(t) + u^{(k)}_t*\Delta t + O(\Delta t)*\Delta t + E'\\
%     &= u^{(k+1)}(t) + u^{(k)}_t*\Delta t + O(\Delta t^2)
% \end{align*}
% Similar to the base case, these local errors will compound for global error $O(\Delta t)$. Combining this with the results from \cite{10.1007/11840817_66} concludes the proof.
% \end{proof}

% \subsection{Proof of Theorem \ref{thm:mprnn_convergence}}
% \label{sec:proof2}
% \begin{proof}
% Let us take a look at the RNN formulation for an individual graph node in MPRNN in Theorem \ref{thm:mprnn_spatiotemporal}. We have
% \begin{align*}
%     \implies u_{p_0,t+\Delta t} &= u_{p_0,t} + i'(t)*\Delta t \\ &+ O((\Delta x + \Delta y + \Delta z + \Delta t)\Delta t + (\Delta t)^2)
% \end{align*}
% Let us consider the case when training samples are provided at the points $p_{0+\Delta x}, p_{0+\Delta y}, p_{0+\Delta z}$. In this case, since the values are not approximated by the RNNs at these points, the corresponding error term changes as follows
% \begin{align*}
%     \implies u_{p_0,t+\Delta t} &= u_{p_0,t} + i'(t)*\Delta t \\ &+ O((\Delta x + \Delta y + \Delta z)\Delta t + (\Delta t)^2)
% \end{align*}
% For the case of uniform spatial grid with $\Delta x = \Delta y = \Delta z$, the above form changes into
% \begin{align*}
%     \implies u_{p_0,t+\Delta t} = u_{p_0,t} + i'(t)*\Delta t + O(\Delta x\Delta t + (\Delta t)^2)
% \end{align*}

% Comparing this form with the definition of linear invariant system from \cite{hardt2018gradient}, we find that the above system is also linear, with a different error model. While the error in \cite{hardt2018gradient} is applied to the output term and assumes a Gaussian distribution, the error distribution for an individual MPRNN node is applied to the input term and depends on the individual partial differential equation and the resolution of the discrete mesh. Based on the definition of the population risk in equation \ref{eqn:population_risk}, we see that $f(\hat{\Theta})$ is symmetric with respect to noise injection in the terms $u_t$ and $\hat{u_t}$. For our specific case, we see that the error is introduced in the $\hat{u_t}$ term, which is $O(\Delta x\Delta t + (\Delta t)^2)$. In fact, the exact form of the error for a single variable for an nth order approximation can be computed as
% \begin{align*}
%     E_n(\Delta t) = \frac{1}{(n+1)!}u^{(n+1)}(t_0+c\Delta t) \Delta t^{n+1}
% \end{align*}
% Generalizing for our case, the error term can be replaced by $Err = V'(u)*(\Delta x\Delta t + (\Delta t)^2)$, where the remaining terms of the series can all be represented by $V'(u)$. Following the proof sketch of \cite{hardt2018gradient}, we find that the variance of the error term appears independently in the additive form of error (see proof of Lemma 3.5 in \cite{hardt2018gradient}). Thus, we can replace the term of variance $\sigma ^2$ in the original result with the variance of our error term.

% Now, for a random variable $e$ and constant $\alpha$, $Var(\alpha e) = \alpha^2 Var(e)$. Since in our error, the terms of $\Delta x$ and $\Delta t$ are constant, we have $Var(Err) = Var(V'(u))*(\Delta x\Delta t + (\Delta t)^2)^2$. We substitute $Var(V'(u))$ with $V(u)$, and replace the $\sigma ^2$ term in the original result to get our result.
% \end{proof}



\section{Problem Setup and Sensor-Space Dynamical Framework}

\paragraph{Physical Models as Dynamical Systems}

We adopt a dynamical systems perspective for physical models, as a wide class of governing equations---including diffusion, advection, and wave phenomena---can be expressed through time-evolution operators acting on latent states. This abstraction allows us to reason about diverse physical processes within a unified framework, independent of the specific form of the underlying PDE.

We consider a physical system defined over a spatial domain $\Omega \subset \mathbb{R}^d$ with state $\mathbf{u}(x,t) \in \mathbb{R}^k$. The system evolves according to a parameterized partial differential equation (PDE):
\begin{equation}
\frac{\partial \mathbf{u}}{\partial t} = \mathcal{F}_\theta(\mathbf{u}, \mathbf{x}, t), \quad \theta \in \Theta,
\end{equation}
where $\mathcal{F}_\theta$ defines the governing dynamics and $\theta$ denotes the set of latent physical parameters.

Given an initial condition $\mathbf{u}(\mathbf{x},0)$, the system induces a trajectory:
\begin{equation}
\mathbf{u}_\theta(\cdot, t) = \Phi_\theta^t(\mathbf{u}(\cdot,0)),
\end{equation}
where $\Phi_\theta^t$ denotes the evolution operator.

\paragraph{Observation Operator}

In real-world sensing applications, the full state of the system is rarely observed. Instead, measurements are obtained through a limited set of sensors that sample the field at discrete spatial locations. The observation operator formalizes this mapping from the latent physical state to measurable quantities.

We define an observation operator $\mathcal{O}$ that maps the full state to sensor measurements. For a set of $S$ sensors located at $\{\mathbf{x}_1, \dots, \mathbf{x}_S\}$:
\begin{equation}
\mathbf{y}_i(t) = \mathbf{u}(\mathbf{x}_i, t), \quad i = 1, \dots, S.
\end{equation}

Thus, the observation operator is given by:
\begin{equation}
\mathbf{y}(t) = \mathcal{O}(\mathbf{u}(\cdot,t)) \in \mathbb{R}^S.
\end{equation}

\paragraph{Longitudinal Sensing Framework}

In many practical deployments, such as environmental monitoring \citep{bhardwaj2025comprehensive, iyer2022modeling} and traffic systems \citep{bhardwaj2023understanding}, sensors are fixed in space and collect measurements over time. This leads to a longitudinal sensing regime, where temporal trajectories at sparse spatial locations provide the primary source of information about the system.

We consider the setting where observations are collected over time, forming a trajectory:
\begin{equation}
\mathbf{Y}_\theta = \{\mathbf{y}(t_1), \mathbf{y}(t_2), \dots, \mathbf{y}(t_T)\} \in \mathbb{R}^{S \times T}.
\end{equation}

This longitudinal sensing regime captures the temporal evolution of the system at sparse spatial locations.

\paragraph{Sensor-Space Dynamical System}

While the underlying dynamics evolve in the high-dimensional state space, observations induce a corresponding evolution in sensor space. Studying this induced system allows us to analyze what aspects of the physical dynamics are preserved or lost under sparse observation.

The evolution of the observed system induces a dynamical relation in sensor space. In particular, let $\mathbf{y}(t)$ denote the sensor observation at time $t$. While the underlying physical system evolves according to
\begin{equation}
\mathbf{u}_\theta(\cdot,t) = \Phi_\theta^t(\mathbf{u}(\cdot,0)), \nonumber
\end{equation}
the observations are given by
\begin{equation}
\mathbf{y}(t) = \mathcal{O}(\mathbf{u}_\theta(\cdot,t)). \nonumber
\end{equation}

Due to the non-invertibility of the observation operator $\mathcal{O}$ under sparse sensing, a given observation $\mathbf{y}(t)$ may correspond to multiple latent states. That is, there may exist distinct $\mathbf{u}_t^{(1)} \neq \mathbf{u}_t^{(2)}$ such that
\begin{equation}
\mathcal{O}(\mathbf{u}_t^{(1)}) = \mathcal{O}(\mathbf{u}_t^{(2)}) = \mathbf{y}(t). \nonumber
\end{equation}
If these latent states evolve to different future observations,
\begin{equation}
\mathcal{O}\!\left(\Phi_\theta(\mathbf{u}_t^{(1)})\right) \neq 
\mathcal{O}\!\left(\Phi_\theta(\mathbf{u}_t^{(2)})\right), \nonumber
\end{equation}
then the same current observation admits multiple valid next observations.

Thus, the evolution in sensor space cannot, in general, be represented as a deterministic Markovian system of the form $\mathbf{y}(t+1) = G_\theta(\mathbf{y}(t))$. Instead, it is more naturally viewed as a set-valued or non-Markovian relation induced by the underlying dynamics. This non-invertibility of the observation operator is the fundamental source of ambiguity in inverse problems under sparse sensing.

\paragraph{Observational Equivalence}
A central question in inverse problems is whether distinct physical configurations can be uniquely distinguished from observational data. In this thread, consider two separate forms of equivalences:
\begin{enumerate}[itemsep=0pt,topsep=0pt]
    \item Parameter non-identifiability arising from different parameter values for the same physics model.
    \item Model Indistinguishability arising from different underlying dynamics.
\end{enumerate}

\textit{\textbf{Parameter Non-Identifiability:}} A central question in inverse problems is whether latent parameters can be uniquely recovered from observations. Under sparse sensing, multiple parameter configurations may induce identical observation trajectories, leading to non-identifiability.

We formalize this by defining an equivalence relation over parameters:
\begin{equation}
\theta_1 \sim \theta_2 \quad \text{if} \quad \mathbf{Y}_{\theta_1} \approx \mathbf{Y}_{\theta_2}.
\end{equation}

This captures the notion of parameter non-identifiability under the given sensing regime, where distinct parameter values give rise to indistinguishable observations.

\textit{\textbf{Model Indistinguishability:}} Beyond parameter non-identifiability within a model class, we are interested in whether different physical models can be distinguished under sparse sensing. Let $\Phi^{(1)}$ and $\Phi^{(2)}$ denote two distinct dynamical systems, potentially corresponding to different governing equations.

We say that the two models are indistinguishable under the observation operator $\mathcal{O}$ if there exist trajectories such that
\begin{equation}
\mathbf{Y}^{(1)} \approx \mathbf{Y}^{(2)},
\end{equation}
where $\mathbf{Y}^{(i)}$ denotes the observation trajectory induced by $\Phi^{(i)}$.

This captures the notion that fundamentally different physical dynamics may produce nearly identical observations under sparse sensing, making them indistinguishable using data alone.







\section{Theoretical Analysis via Linear Dynamical Systems}
\label{sec:theory}

In this section we cast the sparse longitudinal sensing problem in a finite-dimensional linear dynamical systems form. The goal is to make the key mechanisms geometric and operator-theoretic rather than PDE-specific: diffusion acts as a smoothing operator, advection acts as a transport operator, and sparse sensing removes precisely the directions in state space where these models differ.

\subsection{Finite-Horizon Observation and Source Maps}

Let the discretized physical field at time $t$ be $\mathbf{u}_t \in \mathbb{R}^n$, where $n$ is the number of grid cells. We assume a linearized or fixed-coefficient discretization of the form
\begin{equation}
\mathbf{u}_{t+1} = A \mathbf{u}_t + B \mathbf{s},
\qquad
\mathbf{y}_t = O \mathbf{u}_t + \boldsymbol{\varepsilon}_t,
\label{eq:lds_setup}
\end{equation}
where:
\begin{itemize}
    \item $A \in \mathbb{R}^{n \times n}$ is the state transition operator,
    \item $B \in \mathbb{R}^{n \times q}$ maps the source coefficients $\mathbf{s}\in \mathbb{R}^q$ into the field,
    \item $O \in \mathbb{R}^{m \times n}$ is the observation operator selecting or interpolating the $m$ sensor measurements,
    \item $\boldsymbol{\varepsilon}_t$ is sensor noise.
\end{itemize}

We consider a finite observation horizon $T$. Stacking the observations
\begin{equation}
\mathbf{Y}
=
\begin{bmatrix}
\mathbf{y}_0\\
\mathbf{y}_1\\
\vdots\\
\mathbf{y}_T
\end{bmatrix}
\in \mathbb{R}^{m(T+1)},
\end{equation}
gives the standard finite-horizon input--output decomposition
\begin{equation}
\mathbf{Y}
=
\mathcal{O}_T(A)\mathbf{u}_0 + \mathcal{S}_T(A,B)\mathbf{s} + \mathbf{E},
\label{eq:stacked_map}
\end{equation}
where $\mathbf{E}$ stacks the sensor noise and
\begin{equation}
\mathcal{O}_T(A)
=
\begin{bmatrix}
O\\
OA\\
\vdots\\
OA^T
\end{bmatrix}
\in \mathbb{R}^{m(T+1)\times n}
\label{eq:obs_map}
\end{equation}
is the finite-horizon observation map, while
\begin{equation}
\mathcal{S}_T(A,B)
=
\begin{bmatrix}
0\\
OB\\
O(I+A)B\\
\vdots\\
O\big(\sum_{k=0}^{T-1}A^k\big)B
\end{bmatrix}
\in \mathbb{R}^{m(T+1)\times q}
\label{eq:source_map}
\end{equation}
is the finite-horizon source map.

The matrix $\mathcal{O}_T(A)$ determines what latent state directions remain visible after projection through the sensors, and $\mathcal{S}_T(A,B)$ determines what source directions remain visible. This is the linear-systems object underlying both distinguishability and identifiability.

\paragraph{Noise-robust equivalence.}
Let $\sigma>0$ denote a sensor noise scale, and let $\|\cdot\|$ denote any operator norm on stacked trajectories. We will say that two observation trajectories are observationally equivalent if their difference is at most of order $\sigma$ in this norm.


\subsection{Model Operators with Open Boundaries}
\label{subsec:model_ops}

We now specialize the general linear dynamical system
\begin{equation}
\mathbf{u}_{t+1} = A \mathbf{u}_t + B \mathbf{s} \nonumber
\end{equation}
to the three canonical classes relevant to sparse environmental sensing: advection, diffusion, and advection--diffusion. Our goal is to make precise, in operator form, the intuition that advection acts primarily as directional transport, diffusion acts primarily as local smoothing, and advection--diffusion differs from advection through a weak smoothing component that may be unobservable under sparse sensing.

\vspace{0.5em}
\noindent
\textbf{Discrete state representation.}
Let the spatial domain be discretized on a regular grid with $n$ cells, and let $\mathbf{u}_t \in \mathbb{R}^n$ denote the stacked field values at time $t$. We assume a fixed time step $\Delta t$ and spatial mesh width $\Delta x$ for simplicity. Unless otherwise stated, $\mathbf{u}_t$ denotes the deviation from a fixed background or baseline field, so that zero boundary inflow corresponds to zero anomalous pollutant entering the domain.

\vspace{0.5em}
\noindent
\textbf{Open-boundary operator.}
In environmental transport settings such as air pollution, the computational domain is not closed: pollutant mass may leave the modeled region through the boundary. We therefore incorporate an open, absorbing boundary condition into the discrete evolution operators.

Let $\mathcal{B}_{\partial}$ denote the linear boundary operator that implements this open-boundary condition. Informally, $\mathcal{B}_{\partial}$ removes mass transported or diffused outside the computational domain and imposes zero anomalous inflow through inflow boundaries. Thus, in the absence of source terms, the boundary cannot create additional pollutant mass.

For nonnegative fields, this boundary condition is mass non-increasing:
\begin{equation}
\|\mathcal{B}_{\partial}\mathbf{z}\|_1
\le
\|\mathbf{z}\|_1,
\qquad
\mathbf{z}\ge 0.
\nonumber
\end{equation}
The superscript $\partial$ below indicates that the corresponding discrete operator includes this open-boundary treatment.

\paragraph{Advection operator.}
Consider first the one-dimensional constant-velocity advection equation
\begin{equation}
\partial_t u + v \partial_x u = 0,
\qquad v>0. \nonumber
\end{equation}
Using an upwind discretization in the interior of the domain,
\begin{equation}
u_i^{t+1}
=
u_i^t - \frac{v\Delta t}{\Delta x}(u_i^t-u_{i-1}^t). \nonumber
\end{equation}
Defining the Courant number
\begin{equation}
c := \frac{v\Delta t}{\Delta x}, \nonumber
\end{equation}
this can be written as
\begin{equation}
u_i^{t+1} = (1-c)u_i^t + c\,u_{i-1}^t. \nonumber
\end{equation}
At an inflow boundary, the ghost-cell value is set to zero anomalous inflow, while at an outflow boundary, mass transported outside the domain is removed. In matrix form, the open-boundary advection update is
\begin{equation}
\mathbf{u}_{t+1}
=
A_{\mathrm{adv}}^{\partial}\mathbf{u}_t,
\qquad
A_{\mathrm{adv}}^{\partial}
=
\mathcal{B}_{\partial}
\left(I-\Delta t\, v D_x\right),
\label{eq:Aadv_def}
\end{equation}
where $D_x$ is the discrete upwind derivative operator and $\mathcal{B}_{\partial}$ encodes the open-boundary treatment.

The key structural property is that $A_{\mathrm{adv}}^{\partial}$ is sparse, directional, and leaky. Each interior row places most of its mass on the current cell and its upstream neighbor. However, unlike a periodic or closed-domain transport operator, $A_{\mathrm{adv}}^{\partial}$ does not preserve mass indefinitely: pollutant transported through the outflow boundary exits the domain. In the ideal transport limit, $A_{\mathrm{adv}}^{\partial}$ behaves like a one-way shift with loss at the boundary, moving information along the wind direction rather than averaging it locally.

\paragraph{Diffusion operator.}
Now consider the one-dimensional diffusion equation
\begin{equation}
\partial_t u = D \partial_{xx} u,
\qquad D>0. \nonumber
\end{equation}
Using a forward-Euler discretization with the standard second-order stencil in the interior,
\begin{equation}
u_i^{t+1}
=
u_i^t + \frac{D\Delta t}{\Delta x^2}(u_{i+1}^t - 2u_i^t + u_{i-1}^t). \nonumber
\end{equation}
Defining
\begin{equation}
r := \frac{D\Delta t}{\Delta x^2}, \nonumber
\end{equation}
we obtain
\begin{equation}
u_i^{t+1} = (1-2r)u_i^t + r\,u_{i-1}^t + r\,u_{i+1}^t. \nonumber
\end{equation}
With open absorbing boundaries, diffusive flux that leaves the domain is removed and no anomalous mass is introduced from outside. The pure diffusion update is therefore
\begin{equation}
\mathbf{u}_{t+1}
=
A_{\mathrm{diff}}^{\partial}\mathbf{u}_t,
\qquad
A_{\mathrm{diff}}^{\partial}
=
\mathcal{B}_{\partial}
\left(I+\Delta t\,D L\right),
\label{eq:Adiff_def}
\end{equation}
where $L$ is the discrete Laplacian with the open-boundary treatment.

Unlike advection, $A_{\mathrm{diff}}^{\partial}$ is approximately symmetric and averaging in the interior: each row mixes neighboring cells in both directions. Repeated powers of $A_{\mathrm{diff}}^{\partial}$ smooth out high-frequency or localized structure, while the open boundary removes mass that diffuses out of the domain. Thus diffusion acts both as a smoothing mechanism and, under absorbing boundaries, as a dissipative source-free evolution.

\paragraph{Advection--diffusion operator.}
For the constant-coefficient advection--diffusion equation
\begin{equation}
\partial_t u + v\partial_x u = D \partial_{xx}u, \nonumber
\end{equation}
combining the above discretizations with the same open-boundary treatment yields
\begin{equation}
\mathbf{u}_{t+1}
=
A_{\mathrm{ad}}^{\partial}\mathbf{u}_t,
\qquad
A_{\mathrm{ad}}^{\partial}
=
\mathcal{B}_{\partial}
\left(I-\Delta t\,vD_x+\Delta t\,D L\right).
\label{eq:Aad_def}
\end{equation}
Equivalently, we decompose the advection--diffusion operator as
\begin{equation}
A_{\mathrm{ad}}^{\partial}
=
A_{\mathrm{adv}}^{\partial}
+
E_{\mathrm{diff}}^{\partial},
\label{eq:Aad_split}
\end{equation}
where
\begin{equation}
E_{\mathrm{diff}}^{\partial}
:=
\mathcal{B}_{\partial}
\left(\Delta t\,D L\right)
\label{eq:Ediff_def}
\end{equation}
is the diffusive perturbation relative to open-boundary advection.

This decomposition is central to our analysis: under advection-dominated dynamics, the perturbation $E_{\mathrm{diff}}^{\partial}$ is weak relative to transport. Moreover, because the boundary is open and absorbing, source-free perturbations are not retained indefinitely in the domain; they are transported, smoothed, and eventually flushed out. This boundary-induced leakiness will be formalized in the next subsection as an exponential stability or finite-residence-time condition on the source-free operators.


% \subsection{Sparse Sensing in the Advection-Dominated Regime}
\subsection{Structural Regime: Sparse Sensing and Directed Transport}

\paragraph{Assumption A1 (Advection-dominated regime).}
To quantify when diffusion is weak relative to transport, we use the cell P\'eclet number
\begin{equation}
\mathrm{Pe}_{\Delta} := \frac{v\Delta x}{D}.
\label{eq:peclet_cell}
\end{equation}
Large $\mathrm{Pe}_{\Delta}$ corresponds to advection-dominated dynamics. Equivalently, in operator form,
\begin{equation}
\frac{\|A_{\mathrm{diff}}^{\partial}\|}{\|A_{\mathrm{adv}}^{\partial}-I\|}
\;\asymp\;
\frac{D/\Delta x^2}{v/\Delta x}
=
\frac{D}{v\Delta x}
=
\mathrm{Pe}_{\Delta}^{-1}.
\label{eq:operator_ratio}
\end{equation}
Hence $\mathrm{Pe}_{\Delta}\gg 1$ implies that the diffusive contribution is small relative to the transport operator.

We now state the structural sensing assumption from which the later operator bounds are derived.

% \paragraph{Assumption A2 (Sparse isolated sensing).}
% The observation operator $O \in \mathbb{R}^{m\times n}$ samples only a small subset of grid cells, and the minimum sensor spacing
% \begin{equation}
% \Delta_s := \min_{i\neq j}\|\mathbf{x}_i-\mathbf{x}_j\|
% \end{equation}
% is large relative to the local stencil scale. In particular, the sensors do not resolve local curvature or short-range smoothing between neighboring cells.

% \paragraph{Assumption A2 (Sparse sensing with transport non-coverage).}
% The observation operator $O \in \mathbb{R}^{m\times n}$ samples only a small subset of grid cells, with sensor locations
% \[
% X_{\mathrm{sens}} = \{\mathbf{x}_1,\dots,\mathbf{x}_m\}.
% \]
% The minimum sensor spacing
% \begin{equation}
% \Delta_s := \min_{i\neq j}\|\mathbf{x}_i-\mathbf{x}_j\|
% \end{equation}
% is large relative to the local stencil scale, so that sensors do not resolve local curvature or short-range smoothing effects.

% In addition, under the advection dynamics over horizon $T$, each sensor observes only a limited set of source cells. Let $C_j(T)$ denote the transport cone of sensor $j$, and suppose that
% \begin{equation}
% |C_j(T)| \le r_T \quad \text{for all } j=1,\dots,m. \nonumber
% \end{equation}
% We assume an extreme sparsity regime in which the total observable region satisfies
% \begin{equation}
% \left|\bigcup_{j=1}^m C_j(T)\right| \le m r_T < q, \nonumber
% \end{equation}
% where $q$ is the total number of source cells.

% Equivalently, there exist source directions that are transport-invisible over the horizon $T$, i.e.,
% \begin{equation}
% \ker \mathcal{S}_T(A_{\mathrm{adv}}^{\partial}, B) \neq \{0\}. \nonumber
% \end{equation}

% In particular, on an $n\times n$ grid with $q=n^2$ and $r_T = O(n)$, this corresponds to a regime where the number of sensors scales as $m = O(n)$ (or density $m/q = O(1/n)$).

\paragraph{Assumption A2 (Sparse sensor geometry).}
The observation operator \(O\in\mathbb{R}^{m\times n}\) samples a
small subset of grid cells, with sensor locations
\[
X_{\mathrm{sens}}=\{x_1,\ldots,x_m\}.
\]
The sensors are spatially sparse and isolated relative to the local
stencil scale. In particular, if
\[
\Delta_s := \min_{i\neq j}\|x_i-x_j\|
\]
denotes the minimum sensor spacing, then \(\Delta_s\) is large compared
with the grid spacing and the one-step numerical stencil. Thus the
sensor network does not resolve local curvature or short-range
smoothing effects.

\paragraph{Assumption A3 (Dominant directed transport).}
Over the finite observation horizon, the velocity field has a persistent
dominant direction. We write
\[
\mathbf v(x,t)=\bar{\mathbf v}+\delta \mathbf v(x,t),
\]
where \(\bar{\mathbf v}\neq 0\) is the mean monsoon-mode transport
direction and \(\delta \mathbf v(x,t)\) is a bounded perturbation. The
flow does not reverse its dominant direction or create persistent closed
recirculation regions over the horizon. Thus pollutant anomalies are
transported predominantly along downstream characteristics associated
with \(\bar{\mathbf v}\).

\paragraph{Assumption A4 (Observation-relative noise scale).}
\label{ass:noise_scale}
The stacked sensor noise $\mathbf{E}$ over the observation horizon is bounded relative to the observation magnitude. In particular, there exists $\rho \in (0,1)$ such that
\begin{equation}
\|\mathbf{E}\| \le \rho\,\|\mathbf{Y}\|,
\end{equation}
where $\mathbf{Y} = \mathcal{O}_T(A)\mathbf{u}_0$ denotes the noiseless observation trajectory. Equivalently, the effective sensor noise scale may be written as
\begin{equation}
\sigma := \rho\,\|\mathbf{Y}\|.
\end{equation}

\paragraph{Assumption A5 (Open absorbing boundary).}
After subtracting a fixed background field, the source-free dynamics are
evolved on a bounded computational domain using open-boundary operators.
The boundary is absorbing and non-blocking: pollutant anomalies that are
transported or diffused out of the computational domain are removed, and
no anomalous pollutant mass enters through inflow boundaries.

Let \(A^{\partial}\) denote either the source-free advection operator
\(A_{\mathrm{adv}}^{\partial}\) or the source-free advection--diffusion
operator \(A_{\mathrm{ad}}^{\partial}\). The corresponding numerical
scheme is assumed to be positivity preserving and mass non-increasing
for nonnegative anomaly fields:
\begin{equation}
\mathbf u\ge 0
\quad\Longrightarrow\quad
A^{\partial}\mathbf u \ge 0,
\qquad
\|A^{\partial}\mathbf u\|_1
\le
\|\mathbf u\|_1.
\label{eq:open_boundary_mass_nonincrease}
\end{equation}

This assumption captures the physical fact that, in an open
air-pollution domain, the boundary cannot create pollutant anomalies:
mass may leave the modeled region, but no additional anomalous mass is
introduced from outside. 
% The finite residence-time and stability
% properties used below will be derived from this open-boundary condition
% together with the directed-transport geometry.

\subsection{Transport Coverage and Source-Space Ambiguity}

\paragraph{Finite-horizon transport cones.}
For each sensor \(j\), define its upstream finite-horizon transport cone
under open-boundary advection as
\[
\begin{aligned}
\mathcal C_j(T)
:=
\bigl\{
\ell\in\{1,\ldots,q\}:\;&
\exists\, t\in\{0,\ldots,T-1\}
\\
&\text{such that }
O_j (A_{\mathrm{adv}}^{\partial})^t B e_\ell \neq 0
\bigr\}.
\label{eq:transport_cone_def}
\end{aligned}
\]
Here \(e_\ell\) denotes the \(\ell\)-th source basis vector and \(O_j\)
is the row of \(O\) corresponding to sensor \(j\). Thus
\(\mathcal C_j(T)\) is the set of source cells whose emissions can reach
sensor \(j\) within the observation horizon under pure advection.

Let
\begin{equation}
r_{j,T}:=|\mathcal C_j(T)|,
\qquad
r_T:=\max_{1\le j\le m} r_{j,T},
\end{equation}
and define the total transport-covered source set as
\begin{equation}
\mathcal C(T)
:=
\bigcup_{j=1}^m \mathcal C_j(T).
\label{eq:covered_source_set}
\end{equation}
Since the sensor cones may overlap,
\begin{equation}
|\mathcal C(T)|
\le
\sum_{j=1}^m |\mathcal C_j(T)|
=
\sum_{j=1}^m r_{j,T}
\le
m r_T.
\label{eq:cone_union_bound}
\end{equation}

\begin{lemma}[Sparse transport coverage implies source-map nullspace]
\label{lem:sparse_transport_noncoverage}
Let \(\mathcal C(T)\) be the finite-horizon transport-covered source set
defined in \eqref{eq:covered_source_set}. If
\begin{equation}
|\mathcal C(T)| < q,
\label{eq:transport_noncoverage_condition}
\end{equation}
then the pure-advection finite-horizon source map has a nontrivial
nullspace:
\begin{equation}
\ker S_T(A_{\mathrm{adv}}^{\partial},B)\neq \{0\}.
\end{equation}
More precisely,
\begin{equation}
\dim \ker S_T(A_{\mathrm{adv}}^{\partial},B)
\ge
q-|\mathcal C(T)|.
\label{eq:nullspace_dimension_bound}
\end{equation}
In particular, the sufficient counting condition
\begin{equation}
m r_T < q
\label{eq:sufficient_sparse_condition}
\end{equation}
implies the existence of transport-invisible source directions.
\end{lemma}

Detailed proof of lemma provided in Appendix \ref{sec:proof_lemma_sparse_transport_noncoverage}.


\paragraph{Transport-visible subspace.}
Let $V_{\mathrm{tr}}\subset\mathbb{R}^n$ denote the subspace spanned by basis vectors whose grid cells lie in the upstream transport cones of the sensors over the horizon $T$. Equivalently, $V_{\mathrm{tr}}$ collects the initial perturbations that can be carried by advection into at least one sensor location during the observation window.

The next lemma shows that transport is visible on this subspace.

\begin{lemma}[Transport visibility on the covered subspace]
\label{lem:transport_visibility}
Suppose the observation horizon $T$ is long enough that the transport operator can carry information from upstream cells in $V_{\mathrm{tr}}$ to at least one sensor location. Then the restriction of the finite-horizon observation map
\begin{equation}
\mathcal{O}_T(A_{\mathrm{adv}}^{\partial})
=
\begin{bmatrix}
O\\
OA_{\mathrm{adv}}^{\partial}\\
\vdots\\
OA_{\mathrm{adv}}^{{\partial}^T}
\end{bmatrix} \nonumber
\end{equation}
to $V_{\mathrm{tr}}$ has nonzero gain. In particular, there exists $\alpha>0$ such that
\begin{equation}
\inf_{\mathbf{u}\in V_{\mathrm{tr}},\,\|\mathbf{u}\|=1}
\|\mathcal{O}_T(A_{\mathrm{adv}}^{\partial})\mathbf{u}\|
\ge \alpha. 
\label{eq:transport_visibility_derived}
\end{equation}
\end{lemma}
Detailed proof of lemma provided in Appendix \ref{sec:proof_lemma_transport_visibility}.

This is the finite-dimensional version of the statement that sensors see upstream transport contributions when the transport horizon is long enough.

\subsection{Open-Boundary Leakage and Stability}
\label{subsec:open_boundary_stability}

% \begin{lemma}[Finite-time leakage for open-boundary advection]
% \label{lem:finite_time_leakage_advection}
% Consider the source-free one-dimensional advection equation
% \[
% \partial_t u + v\partial_x u = 0,
% \qquad v>0,
% \]
% on a bounded interval of length \(L\), with zero anomalous inflow at
% the upstream boundary and absorbing outflow at the downstream boundary.
% Let the domain be discretized into \(H\) streamwise cells, where
% \[
% H = \left\lceil \frac{L}{\Delta x} \right\rceil,
% \]
% and let the upwind Courant number be
% \[
% c = \frac{v\Delta t}{\Delta x},
% \qquad 0<c\le 1.
% \]
% Let \(A_{\mathrm{adv}}^{\partial}\) denote the resulting open-boundary
% upwind advection operator. Then \(A_{\mathrm{adv}}^{\partial}\) is
% mass non-increasing in the \(\ell_1\) norm:
% \[
% \|A_{\mathrm{adv}}^{\partial}\mathbf{u}\|_1
% \le
% \|\mathbf{u}\|_1,
% \qquad
% \mathbf{u}\ge 0.
% \]
% Moreover, after
% \[
% N_{\mathrm{exit}} = H
% \]
% steps, at least a fraction
% \[
% \eta = c^H
% \]
% of the mass of any nonnegative initial distribution has exited the
% domain. Equivalently,
% \begin{equation}
% \|(A_{\mathrm{adv}}^{\partial})^{N_{\mathrm{exit}}}\mathbf{u}\|_1
% \le
% (1-\eta)\|\mathbf{u}\|_1,
% \qquad
% \mathbf{u}\ge 0.
% \label{eq:adv_leakage_derived}
% \end{equation}
% \end{lemma}

\begin{lemma}[Finite-time leakage for open-boundary advection]
\label{lem:finite_time_leakage_advection}
Consider source-free one-dimensional advection on a bounded interval of
length \(L\), with velocity \(v>0\), zero anomalous inflow, and absorbing
outflow. Let
\[
H=\left\lceil \frac{L}{\Delta x}\right\rceil,
\qquad
c=\frac{v\Delta t}{\Delta x},
\qquad
0<c\le 1,
\]
and let \(A_{\mathrm{adv}}^{\partial}\) be the open-boundary upwind
advection operator. Then \(A_{\mathrm{adv}}^{\partial}\) is
\(\ell_1\)-mass non-increasing for nonnegative fields:
\[
\|A_{\mathrm{adv}}^{\partial}\mathbf u\|_1
\le
\|\mathbf u\|_1,
\qquad
\mathbf u\ge 0.
\]
Moreover, after \(H\) steps, at least a fraction \(c^H\) of any
nonnegative initial mass exits the domain. Equivalently,
\begin{equation}
\|(A_{\mathrm{adv}}^{\partial})^{H}\mathbf u\|_1
\le
(1-c^H)\|\mathbf u\|_1,
\qquad
\mathbf u\ge 0.
\label{eq:adv_leakage_derived}
\end{equation}
\end{lemma}

Detailed proof provided in Appendix \ref{sec:proof_finite_time_leakage_advection}.

\begin{corollary}[Exponential stability of open-boundary advection]
\label{cor:adv_exp_stability_derived}
Under the conditions of Lemma~\ref{lem:finite_time_leakage_advection},
the open-boundary advection operator is exponentially stable in
\(\ell_1\). In particular,
\[
\|(A_{\mathrm{adv}}^{\partial})^t\|_1
\le
M_{\mathrm{adv}}\lambda_{\mathrm{adv}}^t,
\qquad t\ge 0.
\]
% where one valid choice is
% \[
% \lambda_{\mathrm{adv}}
% =
% (1-c^H)^{1/H}
% \]
% and
% \[
% M_{\mathrm{adv}}
% =
% \max_{0\le r<H}
% \|(A_{\mathrm{adv}}^{\partial})^r\|_1
% \lambda_{\mathrm{adv}}^{-r}.
% \]
% Since \(A_{\mathrm{adv}}^{\partial}\) is non-expansive in \(\ell_1\),
% one may take
% \[
% M_{\mathrm{adv}}
% \le
% \lambda_{\mathrm{adv}}^{-(H-1)}.
% \]
\end{corollary}

In multiple dimensions, the same argument applies along the streamwise
coordinate. For fixed wind direction \(\widehat{\mathbf v}\), define
\[
L_{\parallel}
=
\sup_{\mathbf x,\mathbf y\in\Omega}
(\mathbf x-\mathbf y)\cdot \widehat{\mathbf v}.
\]
The continuous source-free residence time is bounded by
\[
T_{\mathrm{exit}}
=
\frac{L_{\parallel}}{\|\mathbf v\|}.
\]
Thus, after
\[
N_{\mathrm{exit}}
=
\left\lceil
\frac{T_{\mathrm{exit}}}{\Delta t}
\right\rceil
\]
time steps, all purely advected pollutant anomalies have reached the
outflow boundary in the ideal transport limit. For the upwind
discretization with Courant number \(c\le 1\), this finite-time
flushing becomes a block leakage condition with a positive leakage
fraction \(\eta>0\), as in Lemma~\ref{lem:finite_time_leakage_advection}.

\paragraph{Open-boundary stability condition.}
The fixed-wind advection operator satisfies the leakage condition by
Lemma~\ref{lem:finite_time_leakage_advection}. We assume the same
source-free leakage condition for the advection--diffusion operator,
which holds for absorbing open-boundary discretizations satisfying
positivity, mass non-increase, and absence of closed trapping regions.


\subsection{Weak Observable Diffusion}
\paragraph{Diffusion contraction on transport-sensitive modes.}
We now show that diffusion does not preserve the same transport-visible structure.

% \begin{lemma}[Diffusion contraction on transport-visible modes]
% \label{lem:diffusion_contraction}
% Under Assumptions A1 and A2, there exists $\beta>0$ such that
% \begin{equation}
% \sup_{\mathbf{u}\in V_{\mathrm{tr}},\,\|\mathbf{u}\|=1}
% \|\mathcal{O}_T(A_{\mathrm{diff}}^{\partial})\mathbf{u}\|
% \le \beta,
% \end{equation}
% where $\beta$ depends on the sensor density and decreases as sensing becomes sparser.
% \end{lemma}

\begin{lemma}[Weak diffusive visibility of separated transport modes]
\label{lem:diffusion_contraction}
Let \(P_{\mathrm{tr}}\) denote the orthogonal projector onto
\(V_{\mathrm{tr}}\), and let \(A_{\mathrm{diff}}^{\partial}\) be the
open-boundary pure-diffusion operator. Define the finite-horizon
diffusion gain on transport-visible modes as
\begin{equation}
\beta_T
:=
\left\|
\mathcal O_T(A_{\mathrm{diff}}^{\partial})
P_{\mathrm{tr}}
\right\|.
\label{eq:beta_T_def}
\end{equation}
Equivalently,
\begin{equation}
\sup_{\mathbf u\in V_{\mathrm{tr}},\,\|\mathbf u\|=1}
\left\|
\mathcal O_T(A_{\mathrm{diff}}^{\partial})\mathbf u
\right\|
=
\beta_T.
\end{equation}

Moreover, suppose the transport-visible cells are separated from the
sensor cells by distance \(d_{\mathrm{sep}}>0\), and the diffusion
length over the observation horizon,
\begin{equation}
\ell_D(T)
:=
\sqrt{D T\Delta t},
\end{equation}
is small relative to this separation:
\begin{equation}
\ell_D(T) \ll d_{\mathrm{sep}}.
\label{eq:diffusion_length_separation}
\end{equation}
Then the diffusion gain \(\beta_T\) is small. In particular, under a
standard discrete heat-kernel bound for the open-boundary diffusion
operator, there exists a constant \(C_D>0\) such that
\begin{equation}
\beta_T
\le
C_D
\exp\left(
-
\frac{
d_{\mathrm{sep}}^2
}{
C_D D T\Delta t
}
\right).
\label{eq:beta_heat_kernel_bound}
\end{equation}
Thus pure diffusion has weak finite-horizon visibility on transport
modes that are separated from the sensor set by more than the diffusion
length.
\end{lemma}

Detailed proof provided in Appendix \ref{sec:proof_diffusion_contraction}.

\paragraph{Weak observable diffusion and noise-robust equivalence.}
We now show that in the advection-dominated regime, the observable difference between advection and advection--diffusion remains small over a finite observation horizon, and may fall below the sensor noise scale.

% \begin{lemma}[Weak observable diffusion under finite-horizon observations]
% \label{lem:weak_observable_diffusion}
% Let
% \begin{equation}
% A_{\mathrm{ad}}^{\partial} = A_{\mathrm{adv}}^{\partial} + A_{\mathrm{diff}}^{\partial}.  \nonumber
% \end{equation}
% Then the finite-horizon observation maps satisfy
% \begin{equation}
% \|\mathcal{O}_T(A_{\mathrm{ad}}^{\partial}) - \mathcal{O}_T(A_{\mathrm{adv}}^{\partial})\|
% \le
% C_T \,\|A_{\mathrm{diff}}^{\partial}\|,
% \label{eq:delta_bound_basic}
% \end{equation}
% where $C_T$ depends only on the horizon $T$ and bounds on $\|A_{\mathrm{adv}}^{\partial}\|$ and $\|A_{\mathrm{ad}}^{\partial}\|$.

% Moreover, in the advection-dominated regime,
% \begin{equation}
% \|A_{\mathrm{diff}}^{\partial}\|
% \;\asymp\;
% \mathrm{Pe}_{\Delta}^{-1}\,\|A_{\mathrm{adv}}^{\partial}-I\|, \nonumber
% \end{equation}
% so that
% \begin{equation}
% \|\mathcal{O}_T(A_{\mathrm{ad}}^{\partial}) - \mathcal{O}_T(A_{\mathrm{adv}}^{\partial})\|
% \le
% C_T\,\mathrm{Pe}_{\Delta}^{-1}.
% \label{eq:delta_bound_pe}
% \end{equation}

% Consequently, if
% \begin{equation}
% C_T\,\mathrm{Pe}_{\Delta}^{-1}\,\|\mathbf{u}_0\|
% \;\le\;
% \sigma,
% \label{eq:noise_condition}
% \end{equation}
% then the observation trajectories generated by $A_{\mathrm{ad}}^{\partial}$ and $A_{\mathrm{adv}}^{\partial}$ are observationally equivalent up to the sensor noise scale.
% Consequently, under Assumption~A3, if
% \begin{equation}
% C_T\,\mathrm{Pe}_{\Delta}^{-1}\,\|\mathbf{u}_0\|
% \le
% \rho\,\|\mathcal{O}_T(A_{\mathrm{ad}}^{\partial})\mathbf{u}_0\|,
% \label{eq:noise_condition_relative}
% \end{equation}
% then the observation trajectories generated by $A_{\mathrm{ad}}^{\partial}$ and $A_{\mathrm{adv}}^{\partial}$ are observationally equivalent up to the sensor noise level.
% \end{lemma}

% \begin{lemma}[Weak observable diffusion under open-boundary stability]
% \label{lem:weak_observable_diffusion}
% Let
% \[
% A_{\mathrm{ad}}^{\partial}
% =
% A_{\mathrm{adv}}^{\partial}
% +
% E_{\mathrm{diff}}^{\partial},
% \]
% where \(E_{\mathrm{diff}}^{\partial}\) is the diffusive perturbation
% relative to open-boundary advection. Suppose the source-free
% open-boundary dynamics satisfy the stability bounds
% \[
% \|(A_{\mathrm{adv}}^{\partial})^t\|
% \le
% M_{\mathrm{adv}}\lambda_{\mathrm{adv}}^t,
% \quad
% \|(A_{\mathrm{ad}}^{\partial})^t\|
% \le
% M_{\mathrm{ad}}\lambda_{\mathrm{ad}}^t,
% \quad
% t\ge 0,
% \]
% for some \(M_{\mathrm{adv}},M_{\mathrm{ad}}<\infty\) and
% \(\lambda_{\mathrm{adv}},\lambda_{\mathrm{ad}}\in(0,1)\). Then, for
% every finite horizon \(T\),
% \begin{equation}
% \left\|
% O_T(A_{\mathrm{ad}}^{\partial})
% -
% O_T(A_{\mathrm{adv}}^{\partial})
% \right\|
% \le
% C_{\infty}
% \left\|
% E_{\mathrm{diff}}^{\partial}
% \right\|,
% \label{eq:uniform_observable_diffusion_bound}
% \end{equation}
% where
% \begin{equation}
% C_{\infty}
% :=
% \|O\|
% \frac{
% M_{\mathrm{ad}}M_{\mathrm{adv}}
% }{
% (1-\lambda_{\mathrm{ad}})(1-\lambda_{\mathrm{adv}})
% }.
% \label{eq:C_infty_def}
% \end{equation}
% In particular, \(C_{\infty}\) is independent of the observation horizon
% \(T\).

% Moreover, in the advection-dominated regime,
% \[
% \left\|
% E_{\mathrm{diff}}^{\partial}
% \right\|
% \lesssim
% \mathrm{Pe}_{\Delta}^{-1},
% \]
% so that
% \begin{equation}
% \left\|
% O_T(A_{\mathrm{ad}}^{\partial})
% -
% O_T(A_{\mathrm{adv}}^{\partial})
% \right\|
% \le
% C_{\infty}\mathrm{Pe}_{\Delta}^{-1}.
% \label{eq:uniform_pe_bound}
% \end{equation}
% Consequently, under the observation-relative noise condition, if
% \begin{equation}
% C_{\infty}\mathrm{Pe}_{\Delta}^{-1}\|\mathbf u_0\|
% \le
% \rho
% \left\|
% O_T(A_{\mathrm{ad}}^{\partial})\mathbf u_0
% \right\|,
% \label{eq:uniform_noise_condition}
% \end{equation}
% then the observation trajectories generated by
% \(A_{\mathrm{ad}}^{\partial}\) and \(A_{\mathrm{adv}}^{\partial}\)
% are observationally equivalent up to the sensor noise level.
% \end{lemma}

\begin{lemma}[Weak observable diffusion under open-boundary stability]
\label{lem:weak_observable_diffusion}
Let
\[
A_{\mathrm{ad}}^{\partial}
=
A_{\mathrm{adv}}^{\partial}
+
E_{\mathrm{diff}}^{\partial},
\]
where \(E_{\mathrm{diff}}^{\partial}\) is the diffusive perturbation
relative to open-boundary advection. Suppose the source-free
open-boundary dynamics satisfy
\[
\|(A_{\mathrm{adv}}^{\partial})^t\|
\le
M_{\mathrm{adv}}\lambda_{\mathrm{adv}}^t,
\qquad
\|(A_{\mathrm{ad}}^{\partial})^t\|
\le
M_{\mathrm{ad}}\lambda_{\mathrm{ad}}^t,
\qquad
t\ge 0,
\]
with \(M_{\mathrm{adv}},M_{\mathrm{ad}}<\infty\) and
\(\lambda_{\mathrm{adv}},\lambda_{\mathrm{ad}}\in(0,1)\).

Then the homogeneous finite-horizon observation maps satisfy
\begin{equation}
\left\|
\mathcal O_T(A_{\mathrm{ad}}^{\partial})
-
\mathcal O_T(A_{\mathrm{adv}}^{\partial})
\right\|
\le
C_{\infty}^{\mathrm{obs}}
\left\|
E_{\mathrm{diff}}^{\partial}
\right\|,
\label{eq:obs_uniform_bound}
\end{equation}
where
\begin{equation}
C_{\infty}^{\mathrm{obs}}
=
\|O\|
\frac{
M_{\mathrm{ad}}M_{\mathrm{adv}}
}{
(1-\lambda_{\mathrm{ad}})(1-\lambda_{\mathrm{adv}})
}.
\label{eq:Cinf_obs_def}
\end{equation}
In particular, \(C_{\infty}^{\mathrm{obs}}\) is independent of \(T\).

Moreover, the forced source-response maps satisfy the per-time bound
\begin{equation}
\max_{0\le t\le T}
\left\|
O\sum_{k=0}^{t-1}
\left[
(A_{\mathrm{ad}}^{\partial})^k
-
(A_{\mathrm{adv}}^{\partial})^k
\right]B
\right\|
\le
C_{\infty}^{\mathrm{src}}
\left\|
E_{\mathrm{diff}}^{\partial}
\right\|,
\label{eq:src_uniform_per_time_bound}
\end{equation}
where
\begin{equation}
C_{\infty}^{\mathrm{src}}
=
\|O\|\|B\|
\frac{
M_{\mathrm{ad}}M_{\mathrm{adv}}
}{
(1-\lambda_{\mathrm{ad}})(1-\lambda_{\mathrm{adv}})
}.
\label{eq:Cinf_src_def}
\end{equation}
Equivalently, if trajectories are measured in a time-normalized norm,
for example
\[
\|Y\|_{\mathrm{rms},T}
=
\left(
\frac{1}{T+1}
\sum_{t=0}^{T}
\|y_t\|^2
\right)^{1/2},
\]
then
\begin{equation}
\left\|
\mathcal S_T(A_{\mathrm{ad}}^{\partial},B)
-
\mathcal S_T(A_{\mathrm{adv}}^{\partial},B)
\right\|_{\mathrm{rms},T}
\le
C_{\infty}^{\mathrm{src}}
\left\|
E_{\mathrm{diff}}^{\partial}
\right\|.
\label{eq:src_uniform_rms_bound}
\end{equation}

In contrast, under the unnormalized stacked Euclidean norm, the forced
source-response bound scales as
\begin{equation}
\left\|
\mathcal S_T(A_{\mathrm{ad}}^{\partial},B)
-
\mathcal S_T(A_{\mathrm{adv}}^{\partial},B)
\right\|_{2}
\le
\sqrt{T+1}\,
C_{\infty}^{\mathrm{src}}
\left\|
E_{\mathrm{diff}}^{\partial}
\right\|.
\label{eq:src_stacked_l2_bound}
\end{equation}
\end{lemma}

Detailed proof provided in Appendix \ref{sec:proof_weak_observable_diffusion}.

Lemmas~\ref{lem:transport_visibility}--\ref{lem:weak_observable_diffusion} give derived versions of the quantities $\alpha$, $\beta$, and the finite-horizon observable perturbation between advection and advection--diffusion. They show that these are not free assumptions, but consequences of two structural conditions: sparse isolated sensing and an advection-dominated transport regime. In particular:
\begin{itemize}[itemsep=0pt,topsep=0pt,leftmargin=*]
    \item transport remains visible because it carries upstream information directly into sensors,
    \item diffusion is weakly observable because its effect is local smoothing rather than directional transport,
    \item advection--diffusion becomes observationally equivalent to advection when the accumulated observable effect of diffusion remains below the sensor noise scale.
\end{itemize}


\subsection{Observational Equivalence Under Sparse Sensing}
\paragraph{Source Non-Identifiability}

We first show that source non-identifiability follows directly from the finite-horizon source map.

\begin{proposition}[Noise-robust source non-identifiability in advection-dominated sparse sensing]
\label{prop:source_nonidentifiable}
Under Assumptions A1, A2, and A3:
\begin{enumerate}[leftmargin=*]
    \item the pure-advection source response is exactly invisible,
    \begin{equation}
    \mathcal{S}_T(A_{\mathrm{adv}}^{\partial},B)\Delta \mathbf{s} = 0;
    \label{eq:adv_source_zero}
    \end{equation}
    \item in the advection-dominated regime, the advection--diffusion source response satisfies
    \begin{equation}
    \|\mathcal{S}_T(A_{\mathrm{ad}}^{\partial},B)\Delta \mathbf{s}\|
    \le \sigma
    \label{eq:source_below_noise_prop}
    \end{equation}
    whenever the observable diffusion perturbation is below the sensor noise floor;
    \item consequently, the source is not identifiable up to the sensor noise level $\sigma$, even when the initial condition is known.
\end{enumerate}
\end{proposition}



\paragraph{Interpretation.}
The proposition makes the mechanism explicit. Under sparse isolated sensing, only source directions that feed into the transport-visible sensor cones leave a strong signature under advection. Source directions outside those cones are exactly invisible under pure advection. In an advection-dominated advection--diffusion regime, diffusion perturbs this invisibility only weakly, and if the perturbation remains below the noise floor, those source directions remain observationally invisible.

\begin{corollary}[Extension to unknown initial conditions]
\label{cor:source_nonid_unknown_ic}
If the source is not identifiable up to noise when the initial condition is known, then it is also not identifiable up to noise when the initial condition is unknown.
\end{corollary}

The unknown-initial-condition problem strictly contains the known-initial-condition problem as a special case. Therefore any non-identifiability present when $\mathbf{u}_0$ is fixed persists when $\mathbf{u}_0$ is also allowed to vary.

\paragraph{Model Indistinguishability}

The next proposition formalizes the fact that weak diffusion becomes invisible after projection through sparse sensors.

\begin{proposition}[Advection and advection--diffusion are $\delta$-indistinguishable]
\label{prop:adv_ad_indist}
By using Lemma \ref{lem:weak_observable_diffusion}:
\begin{equation}
\|\mathcal{O}_T(A_{\mathrm{ad}}^{\partial}) - \mathcal{O}_T(A_{\mathrm{adv}}^{\partial})\| \le \delta.
\end{equation}
Then for every initial state $\mathbf{u}_0$,
\begin{equation}
\|\mathcal{O}_T(A_{\mathrm{ad}}^{\partial})\mathbf{u}_0
-
\mathcal{O}_T(A_{\mathrm{adv}}^{\partial})\mathbf{u}_0\|
\le
\delta \|\mathbf{u}_0\|.
\label{eq:adv_ad_bound}
\end{equation}
In particular, if $\delta \|\mathbf{u}_0\| \le \sigma$, then the advection and advection--diffusion trajectories are observationally indistinguishable up to the sensor noise level.
\end{proposition}



\paragraph{Interpretation.}
The lemma says that the observable difference between advection and advection--diffusion is controlled by the size of the diffusion perturbation $A_{\mathrm{diff}}^{\partial}$ after repeated propagation and projection through $O$. Under extreme sparsity, the rows of $O$ sample only a few locations and therefore suppress precisely the local smoothing modes created by $A_{\mathrm{diff}}^{\partial}$. This is the operator-theoretic version of the statement that diffusion is present in the latent field but does not survive projection into sensor space.


As a corollary, we can say that advection and diffusion are distinguishable model classes.

\begin{corollary}[Advection and diffusion are distinguishable on transport-sensitive states]
\label{prop:adv_diff_dist}
Under Lemmas \ref{lem:transport_visibility} and \ref{lem:diffusion_contraction}, for every unit vector $\mathbf{u}\in V_{\mathrm{tr}}$ we have
\begin{equation}
\|\mathcal{O}_T(A_{\mathrm{adv}}^{\partial})\mathbf{u}
-
\mathcal{O}_T(A_{\mathrm{diff}}^{\partial})\mathbf{u}\|
\ge
\alpha-\beta.
\label{eq:adv_diff_sep}
\end{equation}
Consequently, if $\alpha-\beta > \sigma$, then advection and diffusion are distinguishable above the sensor noise floor on the transport-sensitive subspace $V_{\mathrm{tr}}$.
\end{corollary}

\paragraph{Interpretation.}
The proposition says that advection and diffusion remain distinguishable whenever there exists a subspace of initial states whose transport signatures survive projection by $O$, while diffusion suppresses those same signatures. This is the linear-algebraic version of the statement that transport moves structure across sensors, whereas diffusion merely smooths locally.



\subsection{Geometric Understanding of the Results}

The linear-algebraic form makes the geometry of the problem explicit.

\begin{itemize}[itemsep=0pt,topsep=0pt,leftmargin=*]
    \item The observation operator $O$ is a projection from a high-dimensional field space to a low-dimensional sensor space.
    \item The advection operator $A_{\mathrm{adv}}^{\partial}$ moves information directionally, so its repeated powers $OA_{\mathrm{adv}}^{\partial}^t$ can create distinguishable time-delayed signatures at sparse sensors.
    \item The diffusion operator $A_{\mathrm{diff}}^{\partial}$ is smoothing/contractive, so repeated powers $OA_{\mathrm{diff}}^{\partial}^t$ suppress the very directions that carry sharp transport structure.
    \item The advection--diffusion operator $A_{\mathrm{ad}}^{\partial} = A_{\mathrm{adv}}^{\partial} + A_{\mathrm{diff}}^{\partial}$ differs from advection primarily through a local smoothing perturbation $A_{\mathrm{diff}}^{\partial}$. Under sparse sensing, that perturbation is weak after projection and can fall below the sensor noise floor.
    \item The source map $\mathcal{S}_T(A,B)$ identifies which source directions can actually reach the sensors over the horizon. Under fixed wind and sparse sensors, only source locations lying in downstream sensor cones contribute appreciably. Sources outside those cones lie in nearly invisible directions.
\end{itemize}

This is the finite-dimensional, operator-theoretic version of the main message of the paper: sparse longitudinal sensing does not merely make inference numerically difficult. It changes the geometry of the observation map so that distinct physical explanations collapse after projection through the sensors.





\section{Proofs}
\subsection{Proof of Lemma \ref{lem:sparse_transport_noncoverage}}
\label{sec:proof_lemma_sparse_transport_noncoverage}
\begin{proof}
If \(\ell\notin \mathcal C(T)\), then by definition source cell \(\ell\)
does not reach any sensor within the observation horizon under
open-boundary advection. Therefore,
\[
O(A_{\mathrm{adv}}^{\partial})^t B e_\ell = 0
\qquad
\text{for all } t=0,\ldots,T-1.
\]
Hence the column of the finite-horizon source map
\(S_T(A_{\mathrm{adv}}^{\partial},B)\) corresponding to source cell
\(\ell\) is zero.

Thus any source perturbation supported on the complement
\(\mathcal C(T)^c\) is invisible to the source map. Since
\[
|\mathcal C(T)^c| = q-|\mathcal C(T)|,
\]
the nullspace has dimension at least \(q-|\mathcal C(T)|\). Finally,
because
\[
|\mathcal C(T)|
\le
m r_T,
\]
the condition \(m r_T<q\) implies \(|\mathcal C(T)|<q\), and therefore
\(\ker S_T(A_{\mathrm{adv}}^{\partial},B)\neq\{0\}\).
\end{proof}


\subsection{Proof of Lemma \ref{lem:transport_visibility}}
\label{sec:proof_lemma_transport_visibility}
\paragraph{Proof sketch.}
By construction of $V_{\mathrm{tr}}$, every nonzero $\mathbf{u}\in V_{\mathrm{tr}}$ has support on at least one grid cell that lies upstream of some sensor and can be transported into that sensor within the observation window. Since $A_{\mathrm{adv}}$ acts as a directional shift along the wind field, there exists some time $t\le T$ such that $OA_{\mathrm{adv}}^t \mathbf{u}\neq 0$. Thus the restriction of $\mathcal{O}_T(A_{\mathrm{adv}})$ to $V_{\mathrm{tr}}$ is injective. Since $V_{\mathrm{tr}}$ is finite-dimensional, the minimum singular value of this restricted map is strictly positive, yielding \eqref{eq:transport_visibility_derived}. \hfill $\square$


\subsection{Proof of Lemma \ref{lem:finite_time_leakage_advection}}
\label{sec:proof_finite_time_leakage_advection}
\begin{proof}
The open-boundary upwind scheme is
\[
u_i^{t+1} = (1-c)u_i^t + c u_{i-1}^t,
\]
with zero anomalous inflow at the upstream boundary and removal of mass
that leaves the downstream boundary. Therefore the update is
positivity preserving and mass non-increasing.

The upwind update can be interpreted as follows: during each time step,
a unit of mass either remains in its current cell with weight \(1-c\) or
moves one cell downstream with weight \(c\). Starting from any cell, the
most upstream possible location requires at most \(H\) downstream moves
to leave the domain. Over \(H\) time steps, the event that the mass moves
downstream at every step has weight \(c^H\). Hence at least a fraction
\(c^H\) of the mass exits the domain within \(H\) steps, regardless of
the initial cell. Therefore
\[
\|(A_{\mathrm{adv}}^{\partial})^H\mathbf{u}\|_1
\le
(1-c^H)\|\mathbf{u}\|_1.
\]
\end{proof}


\subsection{Proof of Lemma \ref{lem:diffusion_contraction}}
\label{sec:proof_diffusion_contraction}
% \paragraph{Proof sketch.}
% The operator $A_{\mathrm{diff}}$ is a local averaging operator: each application redistributes mass from a given cell to its neighboring cells. Repeated application of $A_{\mathrm{diff}}$ therefore spreads any initially localized or structured signal over an increasingly large spatial region.

% Let $\mathbf{u}\in V_{\mathrm{tr}}$ with $\|\mathbf{u}\|=1$. By definition of $V_{\mathrm{tr}}$, $\mathbf{u}$ is supported on a set of cells that contribute to sensor observations under advection. However, under diffusion, this localized structure is progressively dispersed across many grid cells.

% Since the observation operator $O$ samples only a small subset of grid locations, the energy of $A_{\mathrm{diff}}^t \mathbf{u}$ captured by the sensors is proportional to the fraction of total mass that remains concentrated at those sampled locations. As diffusion spreads the signal over a growing neighborhood, the magnitude observed at any fixed sensor decreases.

% Formally, for each $t\le T$, the vector $A_{\mathrm{diff}}^t \mathbf{u}$ has support over a region whose size grows with $t$, while its $\ell_2$ norm remains bounded. Therefore, the restriction of this vector to the $m$ sensor coordinates satisfies
% \begin{equation}
% \|O A_{\mathrm{diff}}^t \mathbf{u}\|
% \le
% \sqrt{\frac{m}{n_t}},
% \end{equation}
% where $n_t$ is the effective number of grid cells over which the signal has spread. As sensing becomes sparser, $m/n_t$ decreases, leading to smaller observed magnitude.

% Stacking over time yields
% \begin{equation}
% \|\mathcal{O}_T(A_{\mathrm{diff}})\mathbf{u}\|
% \le \beta,
% \end{equation}
% for some $\beta$ that decreases with sensor sparsity. \hfill $\square$

\paragraph{Proof sketch.}
By definition,
\[
\beta_T
=
\left\|
\mathcal O_T(A_{\mathrm{diff}}^{\partial})
P_{\mathrm{tr}}
\right\|
\]
is the finite-horizon gain of the pure-diffusion dynamics restricted to
the transport-visible subspace.

For a source-free diffusion operator with absorbing open boundaries,
the discrete heat kernel satisfies a Gaussian off-diagonal decay bound:
for grid cells \(i\) and \(j\),
\[
\left|
\left[
(A_{\mathrm{diff}}^{\partial})^t
\right]_{ji}
\right|
\le
C_D
\exp\left(
-
\frac{
\operatorname{dist}(x_i,x_j)^2
}{
C_D D t\Delta t
}
\right),
\qquad t\ge 1.
\]
If every cell in the support of \(V_{\mathrm{tr}}\) is at least
\(d_{\mathrm{sep}}\) away from every sensor cell, then for all
\(t\le T\),
\[
\left\|
O(A_{\mathrm{diff}}^{\partial})^t P_{\mathrm{tr}}
\right\|
\le
C_D
\exp\left(
-
\frac{
d_{\mathrm{sep}}^2
}{
C_D D T\Delta t
}
\right).
\]
Stacking over the finite horizon and absorbing horizon-dependent
polynomial factors into \(C_D\) gives
\[
\left\|
\mathcal O_T(A_{\mathrm{diff}}^{\partial})
P_{\mathrm{tr}}
\right\|
\le
C_D
\exp\left(
-
\frac{
d_{\mathrm{sep}}^2
}{
C_D D T\Delta t
}
\right).
\]
This proves the claim.
\hfill \(\square\)


\subsection{Proof of Lemma \ref{lem:weak_observable_diffusion}}
\label{sec:proof_weak_observable_diffusion}
% \paragraph{Proof sketch.}
% Using the stacked observation map
% \begin{equation}
% \mathcal{O}_T(A)
% =
% \begin{bmatrix}
% O\\
% OA\\
% \vdots\\
% OA^T
% \end{bmatrix}, \nonumber
% \end{equation}
% we have
% \begin{equation}
% \mathcal{O}_T(A_{\mathrm{ad}}) - \mathcal{O}_T(A_{\mathrm{adv}})
% =
% \begin{bmatrix}
% 0\\
% O(A_{\mathrm{ad}} - A_{\mathrm{adv}})\\
% \vdots\\
% O(A_{\mathrm{ad}}^T - A_{\mathrm{adv}}^T)
% \end{bmatrix}. \nonumber
% \end{equation}

% Since $A_{\mathrm{ad}} - A_{\mathrm{adv}} = A_{\mathrm{diff}}$, we use the telescoping expansion
% \begin{equation}
% A_{\mathrm{ad}}^t - A_{\mathrm{adv}}^t
% =
% \sum_{j=0}^{t-1}
% A_{\mathrm{ad}}^{\,t-1-j}
% A_{\mathrm{diff}}
% A_{\mathrm{adv}}^{\,j}. \nonumber
% \end{equation} 
% This is a standard telescoping identity. For $t=1$ it is trivial. Assume it holds for $t$. Then
% \begin{align}
% A_{\mathrm{ad}}^{t+1}-A_{\mathrm{adv}}^{t+1} \nonumber
% &=
% A_{\mathrm{ad}}^{t+1}-A_{\mathrm{ad}}^tA_{\mathrm{adv}} + A_{\mathrm{ad}}^tA_{\mathrm{adv}}-A_{\mathrm{adv}}^{t+1} \nonumber\\
% &=
% A_{\mathrm{ad}}^t(A_{\mathrm{ad}}-A_{\mathrm{adv}})
% +
% (A_{\mathrm{ad}}^t-A_{\mathrm{adv}}^t)A_{\mathrm{adv}} \nonumber \\
% &=
% A_{\mathrm{ad}}^t A_{\mathrm{diff}}
% +
% \left(\sum_{j=0}^{t-1} A_{\mathrm{ad}}^{\,t-1-j} A_{\mathrm{diff}} A_{\mathrm{adv}}^{\,j}\right)A_{\mathrm{adv}} \nonumber \\
% &=
% \sum_{j=0}^{t} A_{\mathrm{ad}}^{\,t-j} A_{\mathrm{diff}} A_{\mathrm{adv}}^{\,j}, \nonumber
% \end{align}
% which proves the identity by induction.
% Taking norms,
% \begin{equation}
% \|A_{\mathrm{ad}}^t - A_{\mathrm{adv}}^t\|
% \le
% \sum_{j=0}^{t-1}
% \|A_{\mathrm{ad}}\|^{t-1-j}
% \|A_{\mathrm{diff}}\|
% \|A_{\mathrm{adv}}\|^{j}. \nonumber
% \end{equation}
% Summing over $t=1,\dots,T$ and factoring out $\|A_{\mathrm{diff}}\|$ yields
% \begin{equation}
% \|\mathcal{O}_T(A_{\mathrm{ad}}) - \mathcal{O}_T(A_{\mathrm{adv}})\|
% \le
% \|O\|\,
% C_T\,
% \|A_{\mathrm{diff}}\|, \nonumber
% \end{equation}
% for a constant $C_T$ depending only on $T$ and operator norms of $A_{\mathrm{adv}}$ and $A_{\mathrm{ad}}$. This gives \eqref{eq:delta_bound_basic}.

% Finally, using the scaling relation
% \begin{equation}
% \|A_{\mathrm{diff}}\|
% \sim
% \frac{D}{v\Delta x}
% =
% \mathrm{Pe}_{\Delta}^{-1}, \nonumber
% \end{equation}
% we obtain \eqref{eq:delta_bound_pe}. The noise condition \eqref{eq:source_noise_condition} then implies observational equivalence. \hfill $\square$

\begin{proof}
We first prove the homogeneous observation-map bound. For \(t\ge 1\),
the telescoping identity gives
\[
(A_{\mathrm{ad}}^{\partial})^t
-
(A_{\mathrm{adv}}^{\partial})^t
=
\sum_{j=0}^{t-1}
(A_{\mathrm{ad}}^{\partial})^{t-1-j}
E_{\mathrm{diff}}^{\partial}
(A_{\mathrm{adv}}^{\partial})^j .
\]
Using the source-free stability bounds,
\begin{align}
\left\|
(A_{\mathrm{ad}}^{\partial})^t
-
(A_{\mathrm{adv}}^{\partial})^t
\right\|
&\le
\sum_{j=0}^{t-1}
\left\|
(A_{\mathrm{ad}}^{\partial})^{t-1-j}
\right\|
\left\|
E_{\mathrm{diff}}^{\partial}
\right\|
\left\|
(A_{\mathrm{adv}}^{\partial})^j
\right\| \nonumber \\
&\le
M_{\mathrm{ad}}M_{\mathrm{adv}}
\left\|
E_{\mathrm{diff}}^{\partial}
\right\|
\sum_{j=0}^{t-1}
\lambda_{\mathrm{ad}}^{t-1-j}
\lambda_{\mathrm{adv}}^j .
\end{align}
Stacking over \(t=1,\ldots,T\), and bounding the finite sum by the
corresponding infinite double series, gives
\begin{align}
\left\|
\mathcal O_T(A_{\mathrm{ad}}^{\partial})
-
\mathcal O_T(A_{\mathrm{adv}}^{\partial})
\right\|
&\le
\|O\|
M_{\mathrm{ad}}M_{\mathrm{adv}}
\left\|
E_{\mathrm{diff}}^{\partial}
\right\|
\sum_{i=0}^{\infty}
\sum_{j=0}^{\infty}
\lambda_{\mathrm{ad}}^i
\lambda_{\mathrm{adv}}^j \nonumber \\
&=
\|O\|
\frac{
M_{\mathrm{ad}}M_{\mathrm{adv}}
}{
(1-\lambda_{\mathrm{ad}})(1-\lambda_{\mathrm{adv}})
}
\left\|
E_{\mathrm{diff}}^{\partial}
\right\|.
\end{align}
This proves \eqref{eq:obs_uniform_bound}.

We now prove the forced source-map bound. The source contribution at
time \(t\) under dynamics \(A\) is
\[
O\sum_{k=0}^{t-1} A^k B\mathbf s.
\]
Therefore the difference between the advection--diffusion and advection
source responses at time \(t\) is
\[
O\sum_{k=0}^{t-1}
\left[
(A_{\mathrm{ad}}^{\partial})^k
-
(A_{\mathrm{adv}}^{\partial})^k
\right]B\mathbf s .
\]
Taking operator norms,
\begin{align}
\left\|
O\sum_{k=0}^{t-1}
\left[
(A_{\mathrm{ad}}^{\partial})^k
-
(A_{\mathrm{adv}}^{\partial})^k
\right]B
\right\|
&\le
\|O\|\|B\|
\sum_{k=0}^{t-1}
\left\|
(A_{\mathrm{ad}}^{\partial})^k
-
(A_{\mathrm{adv}}^{\partial})^k
\right\| .
\end{align}
The \(k=0\) term is zero. For \(k\ge 1\), applying the telescoping bound
above gives
\begin{align}
\sum_{k=1}^{t-1}
\left\|
(A_{\mathrm{ad}}^{\partial})^k
-
(A_{\mathrm{adv}}^{\partial})^k
\right\|
&\le
M_{\mathrm{ad}}M_{\mathrm{adv}}
\left\|
E_{\mathrm{diff}}^{\partial}
\right\|
\sum_{k=1}^{t-1}
\sum_{j=0}^{k-1}
\lambda_{\mathrm{ad}}^{k-1-j}
\lambda_{\mathrm{adv}}^j \nonumber \\
&\le
M_{\mathrm{ad}}M_{\mathrm{adv}}
\left\|
E_{\mathrm{diff}}^{\partial}
\right\|
\sum_{i=0}^{\infty}
\sum_{j=0}^{\infty}
\lambda_{\mathrm{ad}}^i
\lambda_{\mathrm{adv}}^j \nonumber \\
&=
\frac{
M_{\mathrm{ad}}M_{\mathrm{adv}}
}{
(1-\lambda_{\mathrm{ad}})(1-\lambda_{\mathrm{adv}})
}
\left\|
E_{\mathrm{diff}}^{\partial}
\right\|.
\end{align}
Thus, for every \(t\),
\[
\left\|
O\sum_{k=0}^{t-1}
\left[
(A_{\mathrm{ad}}^{\partial})^k
-
(A_{\mathrm{adv}}^{\partial})^k
\right]B
\right\|
\le
C_{\infty}^{\mathrm{src}}
\left\|
E_{\mathrm{diff}}^{\partial}
\right\|,
\]
which proves the per-time source-response bound
\eqref{eq:src_uniform_per_time_bound}.

The time-normalized RMS bound follows immediately from the per-time
bound, since every time block is bounded by the same constant. Under the
unnormalized stacked Euclidean norm, stacking \(T+1\) uniformly bounded
blocks introduces a factor \(\sqrt{T+1}\), giving
\eqref{eq:src_stacked_l2_bound}.
\end{proof}


\subsection{Proof of Proposition \ref{prop:source_nonidentifiable}}

\paragraph{Proof.}
We prove the three claims in order.

\emph{Step 1: exact invisibility under pure advection.}
By assumption A2, the support of $B\Delta \mathbf{s}$ lies outside the downstream transport cones of all sensors over the horizon $T$. Under the pure advection operator $A_{\mathrm{adv}}$, transport is directional and confined to those cones. Hence no component of the propagated source reaches any observed sensor at any time step $t=0,\dots,T$. Therefore,
\begin{equation}
O A_{\mathrm{adv}}^k B \Delta \mathbf{s} = 0
\qquad
\text{for all } k=0,1,\dots,T-1. \nonumber
\label{eq:adv_source_each_step_zero}
\end{equation}
By the definition of the finite-horizon source map,
\begin{equation}
\mathcal{S}_T(A_{\mathrm{adv}},B)
=
\begin{bmatrix}
0\\
OB\\
O(I+A_{\mathrm{adv}})B\\
\vdots\\
O\left(\sum_{k=0}^{T-1}A_{\mathrm{adv}}^k\right)B
\end{bmatrix}, \nonumber
\end{equation}
every block vanishes on $\Delta \mathbf{s}$. Hence
\begin{equation}
\mathcal{S}_T(A_{\mathrm{adv}},B)\Delta \mathbf{s}=0. \nonumber
\end{equation}
This proves \eqref{eq:adv_source_zero}.

\emph{Step 2: weak observable diffusion under advection dominance.}
Now consider the full advection--diffusion dynamics
\begin{equation}
A_{\mathrm{ad}} = A_{\mathrm{adv}} + A_{\mathrm{diff}}. \nonumber
\end{equation}
Since $\Delta \mathbf{s}$ is invisible under pure advection, any nonzero observed source response under $A_{\mathrm{ad}}$ must arise entirely from the diffusive perturbation. Indeed,
\begin{align}
\mathcal{S}_T(A_{\mathrm{ad}},B)\Delta \mathbf{s}
&=
\big(\mathcal{S}_T(A_{\mathrm{ad}},B)-\mathcal{S}_T(A_{\mathrm{adv}},B)\big)\Delta \mathbf{s}. \nonumber
\label{eq:source_map_difference_identity}
\end{align}
Taking norms gives
\begin{equation}
\|\mathcal{S}_T(A_{\mathrm{ad}},B)\Delta \mathbf{s}\|
\le
\|\mathcal{S}_T(A_{\mathrm{ad}},B)-\mathcal{S}_T(A_{\mathrm{adv}},B)\|\,\|\Delta \mathbf{s}\|. \nonumber
\label{eq:source_map_norm_bound}
\end{equation}

The same perturbation argument used in Lemma~\ref{lem:weak_observable_diffusion} for the observation map applies to the source map as well: there exists a finite-horizon constant $C_T$ such that
\begin{equation}
\|\mathcal{S}_T(A_{\mathrm{ad}},B)-\mathcal{S}_T(A_{\mathrm{adv}},B)\|
\le
C_T \,\|A_{\mathrm{diff}}\|. \nonumber
\label{eq:source_map_perturb_bound}
\end{equation}
Moreover, in the advection-dominated regime, Lemma~\ref{lem:weak_observable_diffusion} gives the scaling
\begin{equation}
\|A_{\mathrm{diff}}\|
\asymp
\mathrm{Pe}_{\Delta}^{-1}\,\|A_{\mathrm{adv}}-I\|. \nonumber
\end{equation}
Hence
\begin{equation}
\|\mathcal{S}_T(A_{\mathrm{ad}},B)-\mathcal{S}_T(A_{\mathrm{adv}},B)\|
\le
C_T\,\mathrm{Pe}_{\Delta}^{-1}. \nonumber
\label{eq:source_map_pe_bound}
\end{equation}
If the advection-dominated regime is strong enough that under assumption A3
\begin{equation}
C_T\,\mathrm{Pe}_{\Delta}^{-1}\,\|\Delta \mathbf{s}\|
\le \sigma, \nonumber
\label{eq:source_noise_condition}
\end{equation}
then we have
\begin{equation}
\|\mathcal{S}_T(A_{\mathrm{ad}},B)\Delta \mathbf{s}\|
\le \sigma. \nonumber
\end{equation}
This proves \eqref{eq:source_below_noise_prop}.

\emph{Step 3: non-identifiability up to noise.}
Now assume the initial condition $\mathbf{u}_0$ is fixed and known. The stacked observations satisfy
\begin{equation}
\mathbf{Y} = \mathcal{O}_T(A_{\mathrm{ad}})\mathbf{u}_0 + \mathcal{S}_T(A_{\mathrm{ad}},B)\mathbf{s} + \mathbf{E}. \nonumber
\end{equation}
Consider the two source configurations
\begin{equation}
\mathbf{s}_1 = \mathbf{s},
\qquad
\mathbf{s}_2 = \mathbf{s}+\Delta \mathbf{s}, \nonumber
\end{equation}
with $\Delta \mathbf{s}\neq 0$. Their difference in stacked observations is
\begin{equation}
\mathbf{Y}(\mathbf{s}_2)-\mathbf{Y}(\mathbf{s}_1)
=
\mathcal{S}_T(A_{\mathrm{ad}},B)\Delta \mathbf{s}. \nonumber
\end{equation}
By \eqref{eq:source_below_noise_prop},
\begin{equation}
\|\mathbf{Y}(\mathbf{s}_2)-\mathbf{Y}(\mathbf{s}_1)\|
\le \sigma. \nonumber
\end{equation}
Thus the two distinct source configurations are observationally equivalent up to the sensor noise level. Therefore the source is not identifiable up to noise, even when the initial condition is known. \hfill $\square$


\subsection{Proof of Proposition \ref{prop:adv_ad_indist}}

\paragraph{Proof.}
This is immediate from the definition of operator norm:
\begin{align}
\|\mathcal{O}_T(A_{\mathrm{ad}})\mathbf{u}_0
-
\mathcal{O}_T(A_{\mathrm{adv}})\mathbf{u}_0\|
&=
\|(\mathcal{O}_T(A_{\mathrm{ad}})-\mathcal{O}_T(A_{\mathrm{adv}}))\mathbf{u}_0\| \nonumber \\
&\le
\|\mathcal{O}_T(A_{\mathrm{ad}})-\mathcal{O}_T(A_{\mathrm{adv}})\|\,
\|\mathbf{u}_0\| \nonumber \\
&\le
\delta \|\mathbf{u}_0\| \nonumber.
\end{align}
If $\delta\|\mathbf{u}_0\|\le \sigma$, the difference is below the sensor noise floor and therefore not distinguishable from noisy observations. \hfill $\square$


\subsection{Proof of Corollary \ref{prop:adv_diff_dist}}
\paragraph{Proof.}
Let $\mathbf{u}\in V_{\mathrm{tr}}$ with $\|\mathbf{u}\|=1$. By the reverse triangle inequality,
\begin{align}
\|\mathcal{O}_T(A_{\mathrm{adv}})\mathbf{u}
-
\mathcal{O}_T(A_{\mathrm{diff}})\mathbf{u}\|
&\ge
\|\mathcal{O}_T(A_{\mathrm{adv}})\mathbf{u}\| \nonumber \\
&\quad -
\|\mathcal{O}_T(A_{\mathrm{diff}})\mathbf{u}\| \nonumber.
\end{align}
Taking the infimum lower bound from Lemma \ref{lem:transport_visibility} and the supremum upper bound from Lemma \ref{lem:diffusion_contraction} yields
\begin{equation}
\|\mathcal{O}_T(A_{\mathrm{adv}})\mathbf{u}
-
\mathcal{O}_T(A_{\mathrm{diff}})\mathbf{u}\|
\ge
\alpha-\beta. \nonumber
\end{equation}
If $\alpha-\beta>\sigma$, then the separation exceeds the sensor noise level, hence the two model classes are distinguishable on $V_{\mathrm{tr}}$. \hfill $\square$







\section{Empirical Validation}
\label{sec:empirical_validation}

We empirically validate the theoretical claims in the pollution setting using the same advection--diffusion simulator and calibration pipeline that underlies our inverse-problem experiments. The empirical setup is designed to mirror the assumptions used in Section~4. First, we fix the initial condition and all known physical parameters, and ask whether two different inverse procedures can recover different source maps that explain the same sparse sensor trajectory equally well. This tests source non-identifiability under known initial condition. Second, we reuse the same tuning framework but selectively turn off advection or diffusion and re-optimize the source map for each model class. This allows us to test whether advection--diffusion is empirically hard to distinguish from pure advection under sparse sensing, while remaining distinguishable from pure diffusion.

\subsection{Pollution Simulator and Sensing Setup}
\label{subsec:pollution_setup}

We simulate pollutant transport on a normalized $[0,1]\times[0,1]$ domain corresponding to the New Delhi region using the two-dimensional advection--diffusion--source equation
\begin{align}
\partial_t u + V_x(t)\,\partial_x u + V_y(t)\,\partial_y u
= k\,\Delta u + S_{\mathrm{known}} + S_{\mathrm{unknown}}. \nonumber
\end{align}
The state is discretized on a $40\times 40$ grid, advection is approximated with first-order upwind differences, diffusion with a five-point Laplacian stencil, and time integration uses a two-stage Heun (RK2) update. The simulator enforces the source decomposition
\begin{equation}
S(x,y) = S_{\mathrm{known}}(x,y) + S_{\mathrm{unknown}}(x,y), \nonumber
\end{equation}
where the known source field is loaded from brick kiln and industrial \citep{GUTTIKUNDA2013101}, population-density \citep{ColumbiaUniversity2018}, and traffic intensity \citep{google_maps} maps, cropped from $80\times 80$ to the $40\times 40$ simulation window and normalized. The unknown source is parameterized on a coarse $10\times 10$ grid and then smoothed and upsampled internally to the simulation grid before rollout.

The initial condition is not learned in these experiments. Instead, it is constructed from government monitoring data \citep{cpcb,cpcb_portal} through Universal Kriging. This is important for the interpretation of the experiment: any ambiguity that remains in the recovered source map is not due to uncertainty in the initial condition, but due to the sensing geometry and the induced observation map.

The physical regime is advection dominated. In the simulator, the base wind is set to $(V_x,V_y)=(1.12,0.984)$ and the diffusivity is fixed to $k=3\times 10^{-4}$, with additional diurnal modulation and AR(1) variability in the wind field. These settings are intended to mimic a monsoon-mode transport regime in which directional transport dominates local smoothing, matching Assumption~A1. The rollout horizon is $T=20$ with $N_t=10{,}000$ time steps and $\Delta t = T/(N_t-1) \approx 2\times 10^{-3}$.

The sensor layout is obtained by mapping real regulatory monitor locations onto the simulation grid. Although the source location file contains 33 monitoring sites, one pair collapses to the same cell after discretization, yielding $m=32$ unique sensor locations on the $40\times 40$ grid. Thus the sensor density is only $m/q = 32/1600 = 0.02$, where $q=40\times 40=1600$ is the number of source cells on the fine grid. This is precisely the sparse longitudinal regime of interest. In the language of Assumption~A2, the sensors form a small, spatially separated subset of the domain, so they cannot resolve local curvature or short-range smoothing effects and do not cover the full transport-reachable source space.

The unknown source is first created on a $10\times 10$ coarse grid as a smooth north-south biased field with localized lobes \citep{bhardwaj2025comprehensive}, then scaled so that its total mass after smoothing and upsampling is approximately $25\%$ of the total known-source mass \citep{GUTTIKUNDA2013101}. Sensor observations are sampled from the full simulated field at the 32 sensor locations.

\subsection{Noise Model}
\label{subsec:noise_model}

Observation noise is injected directly into the simulated sensor trajectories. The clean sensor matrix $Y_{\mathrm{clean}} \in \mathbb{R}^{m\times T}$ is perturbed as
\begin{equation}
Y_{\mathrm{noisy}} = Y_{\mathrm{clean}} + \varepsilon,
\qquad
\varepsilon_{ij} \sim \mathcal{N}(0,\sigma^2), \nonumber
\end{equation}
with
\begin{equation}
\sigma = 0.1\,\max_{i,t} |Y_{\mathrm{clean}}(i,t)|. \nonumber
\end{equation}
Thus the noise scale is \emph{relative} to the magnitude of the observed signal, rather than being an absolute constant. This matches the observation-relative noise model assumed in Section~4 and provides the empirical realization of Assumption~A3.

This choice is also consistent with prior work on air-pollution monitoring, which reports substantial label noise and large substitution errors even between nearby regulatory sensors. In particular, prior work notes that pollution labels are affected by unreliable sensing instrumentation and that even sensors within a 2-km radius can exhibit substantial reading differences, with observed substitution errors on the order of tens of $\mu g/m^3$ \citep{bhardwaj2022learning}. We therefore use a relative $10\%$ observation noise model as a controlled simulation counterpart to this empirically noisy sensing regime \citep{bhardwaj2022learning}.

\subsection{Inverse Solvers: SimGrad and STAMP}
\label{subsec:inverse_solvers}

We compare two inverse procedures built on the same differentiable simulator.

\paragraph{SimGrad.}
The first method, SimGrad, directly optimizes the unknown source map by differentiating through the simulator. The optimization variable is the coarse $10\times 10$ unknown source field, represented as a factorized nonnegative amplitude-shape parameterization,
\begin{equation}
S_{\mathrm{unknown}} = a \cdot B, \nonumber
\end{equation}
where both the scalar amplitude $a$ and the spatial base map $B$ are parameterized through shifted softplus transforms to preserve nonnegativity while keeping gradients stable near zero. The learned coarse map is smoothed and upsampled inside the simulator before rollout.

The calibration objective uses the simulated sensor trajectory $\widehat{Y}$ and the observed trajectory $Y$ and combines data fidelity with spatial regularization on the recovered source field. Concretely, the objective includes a sensor-space mean squared error term together with total variation, Laplacian smoothness, $\ell_2$, and box penalties. In addition, the tuner uses several stabilization strategies that matter in this highly underdetermined regime: a time-horizon curriculum that starts from short prefixes of the trajectory and gradually expands to the full horizon, a decaying drift penalty that prevents systematic mismatch in the mean trajectory early in optimization, and batched multi-start calibration with softmin aggregation across candidate solutions. In the default pollution configuration, the tuner uses a batch of 64 candidate source maps, AdamW optimization, and a 70/30 temporal train/validation split.

\paragraph{STAMP.}
The second method, STAMP (Spatio-Temporal Adaptive Message Passing), uses the same differentiable simulator, the same source parameterization, and the same sensor-space data term, but augments the objective with a dynamics-consistency penalty computed by a frozen message-passing recurrent neural network (MPRNN) prior. The added loss is
\begin{equation}
\mathcal{L}_{\mathrm{STAMP}} = \mathcal{L}_{\mathrm{data}} + \lambda_{\mathrm{stamp}}\,\mathcal{L}_{\mathrm{dyn}}, \nonumber
\end{equation}
where $\mathcal{L}_{\mathrm{dyn}}$ is the one-step prediction error of the pretrained MPRNN \citep{iyer2020forecasting,iyer2022modeling} on the \emph{predicted} sensor trajectory after per-sensor normalization using statistics from the observed trajectory.

The frozen prior is loaded from a checkpoint and evaluated on a geometry-based graph over the sensor locations, and the prior consumes edge features derived from inter-sensor geometry. The dynamics term used in the calibrator is teacher-forced one-step MSE over the normalized predicted sensor sequence. To keep the added prior term from overwhelming the simulator gradient, the tuner optionally rebalances $\lambda_{\mathrm{stamp}}$ using the ratio of gradient norms of the base calibration loss and the prior loss with respect to the source parameters.

The MPRNN prior itself is trained separately. The training setup uses a geometry-based sensor graph, per-sensor normalization from the training split, data-driven heteroscedastic weighting based on a robust scale estimate from first differences, and a loss that combines one-step weighted MSE with a multi-step rollout term. This matters here because STAMP does not impose a generic smoothness prior on the source map; instead, it prefers source maps whose induced sensor trajectories are dynamically plausible under a learned sensor-space prior.

\subsection{Evaluation Results}
\label{subsec:evaluation_protocol}

Given a calibrated source map, we rerun the pollution simulator with the recovered source, sample the predicted sensor trajectory at the same sensor locations, and compare both the recovered source field and the sensor trajectory against the ground truth stored in the dataset. The evaluation reports:
\begin{itemize}[itemsep=0pt,topsep=0pt,leftmargin=*]
    \item source-map error: MSE, RMSE, MAE, and NRMSE between the recovered and true $S_{\mathrm{unknown}}$,
    \item sensor-fit error: MSE, RMSE, MAE, and NRMSE between predicted and observed sensor trajectories,
    \item sensor-trajectory correlation: global Pearson correlation on the stacked sensor trajectory and mean temporal correlation per sensor,
    \item source-structure correlation: flat Pearson correlation of the recovered source with the true source and correlation between their spatial gradient fields.
\end{itemize}


\textit{\textbf{Source Non-Identifiability Validation:}} To test source non-identifiability, we keep the initial condition fixed and known, keep the known source field fixed, and keep the advection--diffusion parameters fixed to their generating values. We then run SimGrad and STAMP independently on the same noisy sensor trajectory. This yields two recovered source maps,
\begin{equation}
\widehat{S}_{\mathrm{SimGrad}} \, , \qquad \widehat{S}_{\mathrm{STAMP}}, \nonumber
\end{equation}
that are optimized against the same sensor observations under the same simulator but under different regularization biases.

\begin{table}[t]
\centering
\small
\begin{tabular}{lcc}
\toprule
\textbf{Metric} & \textbf{SimGrad} & \textbf{STAMP} \\
\midrule
\multicolumn{3}{c}{\textit{Sensor Fit (Observed Trajectory)}} \\
RMSE & 0.00996 & 0.01005 \\
MAE  & 0.00778 & 0.00783 \\
NRMSE & 0.00789 & 0.00796 \\
Flat Pearson Corr. & 0.99923 & 0.99924 \\
Mean Temporal Corr. & 0.99958 & 0.99958 \\
\midrule
\multicolumn{3}{c}{\textit{Recovered Source Field Error}} \\
RMSE & 0.02961 & 0.02973 \\
MAE  & 0.02738 & 0.02749 \\
NRMSE & 0.8246 & 0.8279 \\
\midrule
\multicolumn{3}{c}{\textit{Recovered Source Structure}} \\
Flat Pearson Corr. & 0.4111 & $-0.0104$ \\
Gradient Corr. & 0.5926 & 0.2936 \\
\bottomrule
\end{tabular}
\caption{
Comparison of SimGrad and STAMP on pollution source recovery under identical sensing and dynamics. 
Both methods achieve nearly identical sensor reconstruction error and correlation, while producing 
significantly different recovered source structures.
}
\label{tab:pol_non_identifiability}
\end{table}

Both SimGrad and STAMP achieve nearly identical reconstruction accuracy on the observed sensor trajectories, with RMSE on the order of $10^{-2}$ and near-perfect spatio-temporal correlation ($>0.999$). This indicates that both recovered source maps produce indistinguishable sensor-level dynamics under the same forward model.

However, the recovered source fields differ substantially in structure. While SimGrad exhibits moderate correlation with the ground-truth source ($0.41$ flat correlation), STAMP produces a source map that is effectively uncorrelated ($-0.01$), with significantly lower gradient-field alignment. Despite these differences, both solutions explain the observations equally well.

This provides direct empirical evidence of non-identifiability: multiple distinct source configurations lie within the effective nullspace of the sensing operator and produce indistinguishable sensor trajectories under advection-dominated dynamics. This behavior is consistent with Assumption~A2 (transport non-coverage) and demonstrates that, even with known initial conditions and fixed dynamics, source recovery remains fundamentally ill-posed under sparse sensing.

\textit{\textbf{Model Indistinguishability Validation:}} To test the model-class claims, we reuse the same tuning framework but alter the forward model class while allowing the source map to re-optimize. Concretely, starting from the same observed sensor trajectory, we calibrate source maps under three forward models: advection--diffusion, pure advection, and pure diffusion. The pure-advection case is obtained by turning off diffusion, while the pure-diffusion case is obtained by turning off advection. In each case we optimize the source map with the same differentiable calibration pipeline (SimGrad) so that the comparison is not confounded by solver differences.

\begin{table}[t]
\centering
\small
\begin{tabular}{lccc}
\toprule
\textbf{Metric} & \textbf{Adv+Diff} & \textbf{Diff Only} & \textbf{Adv Only} \\
\midrule
\multicolumn{4}{c}{\textit{Sensor Fit (Observed Trajectory)}} \\
RMSE & 0.00996 & 0.08080 & 0.00983 \\
MAE  & 0.00778 & 0.07244 & 0.00772 \\
NRMSE & 0.00789 & 0.06399 & 0.00779 \\
Flat Pearson Corr. & 0.99923 & 0.99974 & 0.99922 \\
Mean Temporal Corr. & 0.99958 & 0.99992 & 0.99957 \\
\bottomrule
\end{tabular}
\caption{
Comparison of calibration performance under different physics regimes. 
Advection-dominated settings (Adv+Diff, Adv Only) yield nearly identical sensor fits, 
while the diffusion-only regime exhibits significantly worse sensor error.
}
\label{tab:pol_regime_comparison}
\end{table}

Despite being calibrated under different forward models, the Advection+Diffusion and Advection Only regimes achieve nearly identical sensor reconstruction error (RMSE $\approx 10^{-2}$) and near-perfect correlation ($>0.999$). This indicates that, after re-optimizing the source map, pure advection is effectively indistinguishable from advection--diffusion in the observed sensor space under sparse sensing. In contrast, the diffusion only regime exhibits an order-of-magnitude higher error, despite maintaining high correlation, indicating a systematic mismatch in amplitude and dynamics that cannot be compensated by source re-optimization. 

This demonstrates that model indistinguishability is governed by the observable transport structure: models that share the same dominant transport pathways (advection-driven) remain indistinguishable after calibration, while models with fundamentally different propagation mechanisms (diffusion-only) become distinguishable even when optimized over the same source space.







% \section{Discussion and Conclusion}

% The results presented in this work demonstrate that non-identifiability and model indistinguishability are not merely theoretical artifacts, but arise naturally in realistic sensing regimes. The assumptions introduced in Section~3—advection-dominated dynamics (A1), sparse sensing with transport non-coverage (A2), and observation-relative noise (A3)—are well aligned with real-world pollution monitoring systems. Urban air quality networks typically consist of a small number of fixed monitoring stations (e.g., 32 sensors over a large metropolitan region), with transport dominated by wind-driven advection and measurements subject to significant noise and variability~\citep{bhardwaj2022learning}. 

% Importantly, these assumptions are not unique to pollution. Similar regimes arise across a wide range of scientific and engineering systems. For example, in oceanography, sparse buoy networks observe large-scale advection of temperature and salinity fields; in atmospheric science, satellite or ground-based measurements sparsely sample evolving weather systems; in traffic flow, fixed sensors capture vehicle density transported along road networks; and in epidemiology, limited testing infrastructure observes disease spread driven by mobility patterns. In all such settings, information propagates along structured transport pathways while observations remain sparse and noisy, leading to inherent ambiguity in reconstructing underlying sources or mechanisms.

% Taken together, these experiments do not merely test whether a particular inverse solver succeeds or fails. Rather, they are designed to show that the ambiguity predicted by the finite-horizon observation geometry is visible in practice: sparse sensing admits multiple source explanations for the same observed trajectory, and it can collapse advection--diffusion onto advection in sensor space while still separating both from pure diffusion.

% These findings have direct implications for both inverse problems and model selection. First, they show that under sparse sensing, recovering fine-grained source structure is fundamentally ill-posed, even when the forward dynamics and initial conditions are known exactly. Second, they demonstrate that model classes themselves may become indistinguishable after calibration, as long as they share similar observable transport behavior. This challenges the common assumption that better optimization or stronger priors alone can resolve such ambiguities.

% Consequently, improving identifiability in these regimes requires changes not only in modeling, but in sensing. Our results suggest a need for alternative sensing paradigms that increase coverage of the transport domain. Examples include mobile sensing platforms (e.g., sensor-equipped vehicles or drones), denser but lower-cost sensor deployments with higher noise tolerance, and hybrid sensing systems that combine fixed and mobile measurements. Such approaches can help break transport-induced non-coverage and improve observability of latent source structure.

% Finally, while this work focuses on advection--diffusion systems, the proposed framework is not limited to this class of dynamics. The same observation-geometry perspective can be extended to other transport-dominated processes, such as wave propagation in shallow water or acoustic systems, where information travels along characteristic fronts rather than diffusive spreads. Extending the analysis to such systems, as well as to more complex multi-scale and nonlinear dynamics, is an important direction for future work.

% In summary, this work highlights a fundamental limitation of sparse sensing in physical systems: when observations do not sufficiently cover the transport domain, both source recovery and model discrimination become intrinsically ambiguous. Addressing this limitation requires rethinking the interplay between sensing, modeling, and inference in large-scale physical systems.











\section{Introduction}

Many real-world spatio-temporal systems, such as environmental processes, fluid transport, and other physical phenomena, are governed by partial differential equations (PDEs) that encode conservation laws and structural dynamics. A general form of such systems is 
\begin{align}
\partial_t \mathbf{u}(\mathbf{x}, t) 
&= \mathcal{F}\big(
    \mathbf{u}(\mathbf{x}, t),\ 
    \nabla \mathbf{u}(\mathbf{x}, t),\ 
    \Delta \mathbf{u}(\mathbf{x}, t),\ 
    \ldots;\ 
    \boldsymbol{\theta} 
\big), \nonumber \\
&\quad \mathbf{x} \in \Omega \subset \mathbb{R}^d,\quad t \in [0, T]
\end{align}
where \( \mathbf{u}(\mathbf{x}, t) \in \mathbb{R}^C \) denotes the evolving physical state and \( \boldsymbol{\theta} \) represents latent parameters such as diffusivity, advection velocity, or source terms. In practice, however, these systems are not fully observed; instead, measurements are obtained from a finite set of sensors, providing sparse and indirect observations of the underlying state.

A central assumption in scientific machine learning is that such sparse measurements are sufficient to infer hidden physical parameters and even distinguish between competing model classes. In this work, we challenge this assumption. We show that under realistic sparse sensing regimes, this problem is fundamentally ill-posed: both source identification and model selection can be impossible, as multiple latent configurations and even distinct physical models can produce indistinguishable observations over finite time horizons.

To formalize this, we consider the discrete-time dynamical system induced by numerical integration of the PDE:
\[
\mathbf{u}_{t+1} = \Phi_\theta(\mathbf{u}_t),
\]
with observations given by a measurement operator
\[
\mathbf{y}_t = \mathcal{O}(\mathbf{u}_t),
\]
where $\mathbf{y}_t$ corresponds to sensor readings at a finite set of spatial locations. This induces a fundamental mismatch: the system evolves in a high-dimensional state space, while observations lie in a much lower-dimensional sensor space. As a result, the mapping from latent parameters and states to observations is many-to-one, leading to intrinsic ambiguity.

Rather than attempting to solve this inverse problem, we study its fundamental limits. We introduce a dynamical systems framework for analyzing the \textbf{observable behavior} of physical models under sparse sensing, focusing on the trajectories induced in sensor space over a finite horizon. This perspective allows us to characterize when different latent configurations or model classes become observationally equivalent.

We formalize two key notions. First, \textbf{identifiability}, where parameters are identifiable if they uniquely determine observations. Second, \textbf{distinguishability}, where two models are distinguishable if they can not induce the same sensor trajectories. Within this framework, we show that under sparse sensing, both properties can fail: source terms are not identifiable, and distinct physical models can become indistinguishable in observation space.

We ground our analysis in the advection--diffusion equation with unknown source terms, motivated by air pollution modeling. In this setting, we show that distinct source configurations can produce identical sensor trajectories, and that advection--diffusion dynamics can become observationally equivalent to pure advection under sparse and spatially isolated sensing. These results demonstrate that ambiguity arises not from algorithmic limitations, but from fundamental constraints imposed by the sensing process.

Our contributions are as follows:
\begin{itemize}[itemsep=0pt,topsep=0pt,leftmargin=*]
    \item We introduce a dynamical systems framework for analyzing the observable behavior of physical models under sparse sensing.
    \item We show that under sparse sensing, source terms are not identifiable: distinct source configurations can produce identical observations.
    \item We show that distinct physical models (e.g., advection and advection--diffusion) can become observationally indistinguishable.
    \item We validate these phenomena empirically in pollution simulations using gradient-based inverse solvers.
\end{itemize}

Our results suggest that the challenge of learning physical models from sparse data is not merely algorithmic, but fundamentally constrained by the information content of observations. This highlights the need for new sensing and modeling strategies that explicitly account for observability limits in physical AI systems.






% \section{Related Work}
% \paragraph{Data Assimilation and Inverse Problems}
% Estimating latent system parameters from partial observations has been extensively studied in inverse problems and data assimilation \citep{tarantola2005inverse, engl1996regularization}. Classical approaches such as Kalman filtering, ensemble Kalman filters \citep{wang2023kalman}, and variational data assimilation methods \citep{carrassi2018data} aim to infer hidden states and parameters from observations. These methods rely on known dynamical models and often require careful tuning of priors, covariance structures, and noise assumptions. However, these approaches typically assume that the underlying system is sufficiently observable, such that latent states and parameters can be inferred from available measurements. In sparse sensing regimes, this assumption breaks down, and multiple parameter configurations may remain consistent with the observed data, leading to fundamental non-identifiability.

% \paragraph{Physics-Informed Machine Learning}
% Physics-informed neural networks (PINNs) and neural operators integrate governing equations into learning frameworks, enabling parameter estimation and system identification from data \citep{raissi2019physics, li2024physics, li2020fourier, bhardwaj2025fieldformer}. These approaches enforce physical constraints through loss functions or architectural priors, allowing for flexible modeling of complex systems. Despite their success, when applied to inverse problems, these approaches implicitly assume that the inverse problem is well-posed, and that the available observations are sufficient to uniquely constrain the underlying parameters. Under sparse sensing, however, this assumption may fail, as distinct parameterizations can produce identical observational trajectories, limiting the effectiveness of purely optimization-based approaches.

% \paragraph{Spatio-Temporal Models}
% Graph neural networks (GNNs), transformers, and sequence models have been widely used to model spatio-temporal processes from sensor data \citep{li2017diffusion, yu2017spatio, wu2019graph, chen2023group, iyer2022modeling}. These models learn representations directly from observations without explicit reliance on governing equations. While these models effectively capture correlations in sensor observations, they operate purely in observation space and do not address whether the underlying physical systems are identifiable from such data. As a result, they cannot distinguish between observationally equivalent physical models.


% \paragraph{Observability and System Identification}
% The question of whether system parameters can be recovered from observations is closely related to classical notions of observability in control theory \citep{chen1984linear}. For linear dynamical systems, observability characterizes whether the internal state can be uniquely determined from output trajectories. Extensions to nonlinear systems and PDEs have been studied in system identification and inverse problem literature \citep{ljung1987theory, banks2012estimation, colton1990inverse}, often under assumptions of sufficient measurement richness. In contrast, we focus on the regime of sparse spatial sensing, where the observation operator induces a low-dimensional projection of the system, and analyze when such projections lead to indistinguishable dynamics across different physical models.


